\documentclass[11pt]{scrartcl}
\usepackage[mdthm,colorsec,noauthor,noasy]{evan}

\DeclareMathOperator{\nul}{null}
\DeclareMathOperator{\range}{range}
\DeclareMathOperator{\spn}{span}
\newcommand{\LL}{\mathcal L}
\newcommand{\PP}{\mathcal P}
\newcommand{\MM}{\mathcal M}
\newcommand{\Fp}{\mathbb F_p}
\newcommand{\traps}{\par\medskip\noindent\textbf{\color{RawSienna}Traps.}\ }
\setlength{\emergencystretch}{3em}

\title{Math 110 Midterm 1 Prep}
\subtitle{Notes Ch.~1--9, HW1--4, practice midterm, SLC packet}
\author{}
\date{}

\ihead{}
\begin{document}
\maketitle

\section{Review}

\subsection{Sets and functions}
Let $f\colon X\to Y$ be a function.
For a subset $U$ of $X$, the \vocab{image} $f(U)$ is the set of all outputs $f(x)$ with $x$ in $U$.
For a subset $V$ of $Y$, the \vocab{preimage} $f^{-1}(V)$ is the set of all inputs $x$ whose output $f(x)$ lands in $V$.
The preimage makes sense for every function. Writing $f^{-1}(V)$ does not mean $f$ has an inverse.

\begin{proposition*}[HW1 P1]
  Going to the preimage and back never gives more than you started with: $f(f^{-1}(V))\subseteq V$.
  Going to the image and back never gives less: $U\subseteq f^{-1}(f(U))$.
  \begin{itemize}
    \ii The first is always an equality exactly when $f$ is surjective.
    \ii The second is always an equality exactly when $f$ is injective.
  \end{itemize}
\end{proposition*}

\traps Do not assume you get back what you started with. Take $f(x)=x^2$ on $\RR$.
With $U=[0,1]$: $f(U)=[0,1]$ and $f^{-1}([0,1])=[-1,1]$, bigger than $U$.
With $V=[-1,1]$: $f^{-1}(V)=[-1,1]$ and $f([-1,1])=[0,1]$, smaller than $V$; the negative part of $V$ is lost because $x^2$ never reaches it.
In general $f(f^{-1}(V))=V\cap\range f$.

\subsection{Fields}
A \vocab{field} is a set of numbers where you can add, subtract, multiply, and divide by anything except $0$, with the usual rules of algebra.
Examples: $\QQ$, $\RR$, $\CC$, and $\Fp=\{0,1,\dots,p-1\}$ with arithmetic mod $p$, when $p$ is prime.

\begin{proposition*}[HW1 P3]
  In a field, if $ab=0$ then $a=0$ or $b=0$.
\end{proposition*}
\begin{proof}
  Suppose $a\ne0$. Then $a$ has an inverse. Multiply both sides of $ab=0$ by $a^{-1}$ to get $b=0$.
\end{proof}
For example, arithmetic mod $6$ is not a field, because $2\cdot3=6=0$ but neither $2$ nor $3$ is $0$.

\begin{proposition*}[Prop.~3, complex conjugation]
  For $w,z\in\CC$: $\overline{w+z}=\bar w+\bar z$, $\overline{wz}=\bar w\,\bar z$, $\bar z=z$ exactly when $z$ is real, and $z\bar z=|z|^2$ is real and $\ge0$.
  So $z^{-1}=\bar z/|z|^2$ for $z\ne0$.
\end{proposition*}

In $\FF_2$ we have $1+1=0$. So $(x+1)^2=x^2+2x+1=x^2+1$, and squaring twice more gives $(x+1)^8=x^8+1$ (HW1 P7).

\traps Any step where you divide by $2$ breaks in $\FF_2$, because $2=0$ there.
Over a finite field, two different polynomials can give the same function. For example, $x^2$ and $x$ give the same output at both points of $\FF_2$.

\subsection{Vector spaces and subspaces}
Examples of vector spaces: $\FF^n$; $\FF^\infty$ (infinite sequences); functions from a set $S$ to $\FF$; $\PP_n(\FF)$ (polynomials of degree at most $n$); $\FF^{m,n}$ ($m\times n$ matrices); $\LL(V,W)$ (linear maps from $V$ to $W$).
Basic facts that follow from the axioms: $0v=0$, $a0=0$, $(-1)v=-v$, and if $av=0$ then $a=0$ or $v=0$ (HW1 P4).

\begin{definition*}
  A \vocab{subspace} of $V$ is a subset $U$ that contains $0$, and is closed under addition and scalar multiplication.
  A quick way to check both closure conditions at once: $\lambda u+w$ is in $U$ whenever $u,w$ are in $U$.

  The \vocab{sum} $U+W$ is the set of all vectors $u+w$ with $u$ in $U$ and $w$ in $W$.
\end{definition*}

\begin{proposition*}[Prop.~5, 6]
  $U\cap W$ is a subspace. $U+W$ is a subspace, and it is the smallest subspace containing $U\cup W$: any subspace that contains both $U$ and $W$ contains all sums $u+w$.
\end{proposition*}
\begin{proof}
  Intersection: $0$ is in both; if $v,w$ are in both, so is $\lambda v+w$. (SLC problem 7.)
  Sum: $\lambda(u_1+w_1)+(u_2+w_2)=(\lambda u_1+u_2)+(\lambda w_1+w_2)$.
  Smallest: $U\cup W\subseteq U+W$ since $u=u+0$ and $w=0+w$; and any subspace containing $U$ and $W$ is closed under addition, so it contains every $u+w$.
\end{proof}

\begin{proposition*}[Prop.~7 / LADR 1.45, what ``direct'' means]
  The sum $U_1+\dots+U_m$ is \vocab{direct}, written $U_1\oplus\dots\oplus U_m$, if every vector in it can be written as $u_1+\dots+u_m$ (with $u_i\in U_i$) in exactly one way.
  Equivalent test: the only way to write $0=u_1+\dots+u_m$ is with every $u_i=0$.
  For two subspaces, the test is simply $U\cap W=\{0\}$.
\end{proposition*}
\begin{proof}
  Uniqueness for every vector implies uniqueness for $0$. Conversely, if $u_1+\dots+u_m=u_1'+\dots+u_m'$, then $\sum(u_i-u_i')=0$, so each $u_i=u_i'$.
  For two subspaces: if $u+w=0$ then $u=-w$ lies in both $U$ and $W$; and if $v\ne0$ lies in both, then $v+(-v)=0$ is a nontrivial way to write $0$.
\end{proof}
Directness says nothing about whether the sum is all of $V$. When you must \emph{use} directness, use unique decomposition. When you must \emph{prove} it, prove the zero test; for two subspaces, prove $U\cap W=\{0\}$.

\begin{proposition*}[Proving $V=U\oplus W$: two of three]
  Let $V$ be finite-dimensional. $V=U\oplus W$ holds as soon as any two of the following hold:
  \begin{enumerate}
    \ii $V=U+W$;
    \ii $U\cap W=\{0\}$;
    \ii $\dim V=\dim U+\dim W$.
  \end{enumerate}
\end{proposition*}
\begin{proof}
  (1)+(2) is the definition. By the dimension formula $\dim(U+W)=\dim U+\dim W-\dim(U\cap W)$: (2)+(3) gives $\dim(U+W)=\dim V$, so $U+W=V$; and (1)+(3) gives $\dim(U\cap W)=0$.
\end{proof}
Default to (2)+(3): the intersection is usually one line and the dimensions are a count. Use (1)+(3) when a splitting formula hands you the sum, such as $A=\tfrac12(A+A^T)+\tfrac12(A-A^T)$. This rule is for two subspaces only.

\begin{proposition*}[Prop.~20, defining a map on a direct sum]
  If $V=U_1\oplus\dots\oplus U_m$ and $T_i\colon U_i\to W$ are linear, there is exactly one linear $T\colon V\to W$ that agrees with $T_i$ on each $U_i$: $T(u_1+\dots+u_m)=T_1u_1+\dots+T_mu_m$.
  It is well defined because the decomposition is unique.
\end{proposition*}

\traps With three or more subspaces, checking that each pair meets only at $0$ is not enough. Three different lines through the origin in $\RR^2$ pass that check, but their sum is not direct.

The union $U\cup W$ is almost never a subspace. It is one only when one of $U$, $W$ contains the other.

Over $\RR$, a set that is nonempty and closed under addition need not be a subspace (the integer points $\ZZ^n$). Over $\Fp$ it always is, because every scalar is $1+1+\dots+1$.

\subsection{Span, independence, bases}
\begin{definition*}
  The \vocab{span} of $v_1,\dots,v_m$ is the set of all linear combinations $a_1v_1+\dots+a_mv_m$.

  The list is \vocab{linearly independent} if the only combination equal to $0$ is the one with all coefficients $0$.

  A \vocab{basis} is a list that is independent and spans $V$. Equivalently, every vector in $V$ can be written as a combination of the list in exactly one way.
\end{definition*}

\begin{proposition*}[Linear dependence lemma, LADR 2.19]
  If a list is dependent, then one of its vectors is a combination of the earlier ones, and you can delete it without changing the span.
\end{proposition*}

\begin{theorem*}[Replacement lemma, Prop.~10 / LADR 2.22]
  Let $e_1,\dots,e_m$ be linearly independent and let $f_1,\dots,f_n$ span $V$. Then $m\le n$.
  Moreover, the $e$'s can be swapped in for $m$ of the $f$'s one at a time, and the list $e_1,\dots,e_m,f_{m+1},\dots,f_n$ (after relabelling the $f$'s) still spans $V$.
\end{theorem*}
\begin{proof}[Proof idea: one in, one out]
  Start with the spanning list $f_1,\dots,f_n$. Write $e_1$ as a combination of the $f$'s; some coefficient is nonzero, so solve for that $f$ and replace it by $e_1$. The list still spans.
  Now write $e_2$ in terms of the current list $e_1,f_2,\dots,f_n$. Independence of the $e$'s says the expression must use at least one $f$ (otherwise $e_2\in\spn(e_1)$), so swap that $f$ out for $e_2$.
  Repeat. Each step uses up one $f$, so if $m>n$ you run out of $f$'s while $e$'s remain, and the leftover $e$ would be a combination of the earlier $e$'s. So $m\le n$.
\end{proof}
\textbf{How the inequality $m\le n$ is used.} It is the reason dimension is well defined, and it is the one-line justification behind all of these:
\begin{itemize}
  \ii Any two bases have the same length: each is independent and the other spans, so each length is $\le$ the other.
  \ii In an $n$-dimensional space, no independent list has more than $n$ vectors, and no spanning list has fewer than $n$.
  \ii ``$n+1$ vectors in $\FF^n$ are dependent'' and ``$n-1$ vectors cannot span $\FF^n$'' need no calculation.
  \ii A subspace of a finite-dimensional space is finite-dimensional with $\dim U\le\dim V$: an independent list in $U$ is independent in $V$, so it has at most $\dim V$ vectors.
  \ii Basis extension: the second half of the lemma says the leftover $f$'s complete $e_1,\dots,e_m$ to a spanning list, and deleting redundant $f$'s gives a basis containing the $e$'s (Prop.~11).
  \ii Any ``impossible map'' argument: if $T\colon V\to W$ were injective with $\dim V>\dim W$, the image of a basis of $V$ would be $\dim V$ independent vectors in $W$, more than a spanning list of $W$ has.
\end{itemize}

\begin{proposition*}[Prop.~9, 11 / LADR 2.30, 2.32]
  You can delete vectors from any spanning list to get a basis (delete a redundant vector, repeat; the linear dependence lemma says one exists while the list is dependent).
  In a finite-dimensional space, you can add vectors to any independent list to get a basis, and the added vectors can be taken from any spanning list you choose.
\end{proposition*}

\begin{example*}[Ex.~18, a basis by degree comparison]
  For any $a\in\RR$, the list $1,\ x-a,\ (x-a)^2,\ \dots,\ (x-a)^n$ is a basis of $\PP_n(\RR)$.
  Independence: the degrees are all different, so in a relation the top-degree term cannot cancel, so its coefficient is $0$; repeat.
  Then count: $n+1$ independent vectors in an $(n+1)$-dimensional space. No need to prove spanning.
  The same move shows that any list of nonzero polynomials with distinct degrees is independent.
\end{example*}

Infinite-dimensional spaces have bases too, defined the same way with finite combinations: $\{1,x,x^2,\dots\}$ is a basis of $\PP(\RR)$, because every polynomial is a finite combination of monomials in exactly one way (Def.~9, Ex.~12).

\traps A span only contains \emph{finite} combinations. So $e_1,e_2,e_3,\dots$ do not span $\FF^\infty$: the sequence $(1,1,1,\dots)$ would need infinitely many terms (HW2 P2).

Whether a list is independent can depend on the field. $(1,1,0),(0,1,1),(1,0,1)$ is a basis of $\RR^3$, but in $\FF_2^3$ the three vectors add up to $0$.

\subsection{Dimension}
The \vocab{dimension} of $V$ is the number of vectors in any basis.
Examples: $\dim\FF^n=n$, $\dim\PP_n=n+1$ (basis $1,x,\dots,x^n$), $\dim\FF^{m,n}=mn$.

\begin{proposition*}[LADR 2.38, 2.39, 2.42]
  Suppose $\dim V=n$.
  \begin{itemize}
    \ii Any $n$ independent vectors form a basis. Any $n$ vectors that span form a basis. So you only need to check one of the two conditions.
    \ii A subspace has dimension at most $n$, and it has dimension exactly $n$ only if it is all of $V$.
  \end{itemize}
\end{proposition*}

\begin{theorem*}[Prop.~15, 16 / LADR 2.43]
  Let $U,W$ be finite-dimensional subspaces of $V$. Then $\dim(U+W)=\dim U+\dim W-\dim(U\cap W)$.

  More precisely: take a basis $v_1,\dots,v_r$ of $U\cap W$, extend it to a basis $v_1,\dots,v_r,u_1,\dots,u_k$ of $U$ and to a basis $v_1,\dots,v_r,w_1,\dots,w_l$ of $W$. Then
  \[ v_1,\dots,v_r,\ u_1,\dots,u_k,\ w_1,\dots,w_l \]
  is a basis of $U+W$, of length $r+k+l=\dim U+\dim W-r$.
\end{theorem*}
\begin{proof}
  The list spans $U+W$ because the first two groups span $U$ and the first and third span $W$.
  Independence: if $\sum\lambda_iv_i+\sum\mu_ju_j+\sum\nu_hw_h=0$, then $\sum\nu_hw_h=-\sum\lambda_iv_i-\sum\mu_ju_j$ is in $U$ and in $W$, so it is in $U\cap W=\spn(v_1,\dots,v_r)$. Since $v_1,\dots,v_r,w_1,\dots,w_l$ is independent, all $\nu_h=0$. Then the rest of the relation is in the basis of $U$, so all $\lambda_i,\mu_j=0$.
\end{proof}
The proof pattern (start from a basis of the intersection, extend twice, show the junk is in the intersection) is worth knowing; it reappears in the proof of rank--nullity.

Consequences:
\begin{itemize}
  \ii $U+W$ is direct exactly when $\dim(U+W)=\dim U+\dim W$.
  \ii The intersection cannot be too small: if $V$ is finite-dimensional, $\dim(U\cap W)\ge\dim U+\dim W-\dim V$.
  \ii (Prop.~17) If the sum $U_1+\dots+U_m$ is direct, then $\dim(U_1\oplus\dots\oplus U_m)=\dim U_1+\dots+\dim U_m$. Concatenating bases of the $U_i$ gives a basis of the sum.
\end{itemize}

\traps To show two subspaces are equal, show that one is inside the other and that they have the same dimension. Having the same dimension alone is not enough.

\subsection{Linear maps, null space, range}
\begin{definition*}
  $T\colon V\to W$ is \vocab{linear} if $T(\lambda v+w)=\lambda Tv+Tw$ for all vectors and scalars.

  The \vocab{null space} $\nul T$ is the set of vectors that $T$ sends to $0$.
  The \vocab{range} $\range T$ is the set of all outputs of $T$.
  Both are subspaces.
\end{definition*}

\begin{proposition*}[Prop.~18, 19 / LADR 3.6--3.8]
  $\LL(V,W)$, the set of linear maps $V\to W$, is a vector space under $(\mu S+T)(v)=\mu Sv+Tv$.
  The composition $ST$ of linear maps is linear. Composition is associative and distributes over addition, and $TI=IT=T$; but $ST\ne TS$ in general.
\end{proposition*}

\begin{proposition*}[LADR 3.4]
  A linear map is decided by what it does to a basis.
  If $v_1,\dots,v_n$ is a basis of $V$, you can pick $w_1,\dots,w_n$ in $W$ freely, and there is exactly one linear map with $Tv_1=w_1,\dots,Tv_n=w_n$.
\end{proposition*}
This is how you build a linear map in a proof: pick a convenient basis and say where each basis vector goes.

\begin{proposition*}[LADR 3.15]
  $T$ is injective exactly when its null space is just $\{0\}$.
\end{proposition*}
\begin{proof}
  If $T$ is injective, then $Tv=0=T0$ forces $v=0$, so $\nul T=\{0\}$. Conversely, suppose $\nul T=\{0\}$ and $Tv=Tw$. Then $T(v-w)=Tv-Tw=0$, so $v-w$ is in the null space. So $v-w=0$, meaning $v=w$.
\end{proof}

\begin{theorem*}[Rank--nullity, LADR 3.21]
  If $V$ is finite-dimensional, then $\dim V=\dim\nul T+\dim\range T$.
\end{theorem*}
\begin{proof}
  Start with a basis $u_1,\dots,u_m$ of the null space, and add vectors $v_1,\dots,v_k$ to make a basis of $V$.
  We show that $Tv_1,\dots,Tv_k$ is a basis of the range. Then $\dim V=m+k=\dim\nul T+\dim\range T$.
  \begin{itemize}
    \ii They span the range: $T$ sends each $u_i$ to $0$, so every output $Tv$ is a combination of $Tv_1,\dots,Tv_k$.
    \ii They are independent: if $\sum a_iTv_i=0$, then $T(\sum a_iv_i)=0$, so $\sum a_iv_i$ is in the null space. That makes it a combination of the $u_j$. Since the $u$'s and $v$'s together are independent, all $a_i=0$.
  \end{itemize}
\end{proof}
Many hard proofs use this same basis: a basis of the null space, extended to a basis of $V$.

\begin{corollary*}
  If $V$ is bigger than $W$ ($\dim V>\dim W$), no linear map $V\to W$ is injective.
  If $V$ is smaller than $W$, no linear map $V\to W$ is surjective.
\end{corollary*}

\begin{proposition*}[Preimage formula, HW3 P1]
  For $V$ finite-dimensional and a subspace $U$ of $W$, the set of vectors that $T$ sends into $U$ has dimension
  \[ \dim T^{-1}(U)=\dim\nul T+\dim(U\cap\range T). \]
\end{proposition*}
\begin{proof}
  Apply rank--nullity to $T$, using only inputs from $T^{-1}(U)$.
  The vectors it sends to $0$ are still $\nul T$.
  Its outputs are exactly the vectors of $U$ that $T$ can reach, which is $U\cap\range T$.
\end{proof}

\begin{proposition*}[HW3 P2]
  If $T^2=0$, then every output of $T$ is sent to $0$ by $T$, so $\range T\subseteq\nul T$.

  If $V$ is finite-dimensional and $\nul T=\range T$, then rank--nullity gives $\dim V=2\dim\range T$, so $\dim V$ must be even.
\end{proposition*}

\traps Rank--nullity only says the dimensions add up. It does not say $V$ splits as $\nul T\oplus\range T$. For $N(x,y)=(y,0)$, the null space and the range are the same line.

A linear map can send an independent list to a dependent one (unless it is injective).
It sends a basis to a list that spans the range, but that list might not be a basis.

\subsection{Matrices}
\begin{definition*}
  Pick a basis $v_1,\dots,v_n$ of $V$ and a basis $w_1,\dots,w_m$ of $W$.
  The matrix $\MM(T)$ has $m$ rows and $n$ columns.
  Its $j$th column lists the coefficients you need to write $Tv_j$ in terms of $w_1,\dots,w_m$.
\end{definition*}

\begin{proposition*}[LADR 3.43, 3.76, Prop.~27]
  Multiplying matrices is the same as composing maps: $\MM(ST)=\MM(S)\MM(T)$.
  Multiplying a matrix by a coordinate vector is the same as applying the map: $\MM(Tv)=\MM(T)\MM(v)$.
  Because of this, matrix multiplication is associative, distributes over addition, and has the identity matrix as identity. It is not commutative.
  The entry formula is $(AB)_{ik}=\sum_j a_{ij}b_{jk}$.
\end{proposition*}

\begin{proposition*}
  $Ax$ is $x_1(\text{column }1)+\dots+x_n(\text{column }n)$.
  So the range of $A$ is the span of its columns, and the rank of $A$ is the dimension of that span.

  The solutions of $Ax=b$, if there are any, are one solution $x_0$ plus anything in the null space.
  Over $\FF_2$, if there is a solution, there are exactly $2^{\,n-\operatorname{rank}A}$ of them, where $n$ is the number of unknowns (columns).
\end{proposition*}

\traps $\MM(T)$ has $\dim W$ rows and $\dim V$ columns, not the other way round.
The matrix changes if you change the bases.

\subsection{Isomorphisms}
\begin{definition*}
  A linear map $T\colon V\to W$ is an \vocab{isomorphism} (or \vocab{invertible}) if there is a linear map $S\colon W\to V$ that undoes it from both sides: $ST=I$ on $V$ and $TS=I$ on $W$.
  When an isomorphism exists, we write $V\cong W$ and say $V$ and $W$ are isomorphic.
\end{definition*}

\begin{proposition*}[LADR 3.60]
  A map has at most one inverse.
  In fact, if $S_1T=I$ and $TS_2=I$, then $S_1=S_2$.
\end{proposition*}
\begin{proof}
  $S_1=S_1(TS_2)=(S_1T)S_2=S_2$.
\end{proof}

\begin{theorem*}[LADR 3.63]
  $T$ is an isomorphism exactly when it is both injective and surjective.
  You never need to check separately that the inverse is linear. That is automatic.
\end{theorem*}
\begin{proof}
  If $T$ has an inverse: when $Tv=0$, apply $T^{-1}$ to get $v=0$, so $T$ is injective. Every $w$ equals $T(T^{-1}w)$, so $T$ is surjective.

  If $T$ is injective and surjective: each $w$ comes from exactly one $v$. Define $T^{-1}w$ to be that $v$.
  To see it is linear: $T$ sends $\lambda T^{-1}w_1+T^{-1}w_2$ to $\lambda w_1+w_2$. Only one vector goes to $\lambda w_1+w_2$, so $T^{-1}(\lambda w_1+w_2)=\lambda T^{-1}w_1+T^{-1}w_2$.
\end{proof}

\begin{theorem*}[LADR 3.70]
  Two finite-dimensional spaces are isomorphic exactly when they have the same dimension.
\end{theorem*}
\begin{proof}
  If $T$ is an isomorphism, its null space is $\{0\}$ and its range is all of $W$. Rank--nullity gives $\dim V=0+\dim W$.
  If the dimensions are equal, pick a basis of each and let $T$ send the first basis to the second (LADR 3.4). $T$ is injective because the image of a relation is a relation among the second basis, and surjective because the second basis spans; so $T$ is an isomorphism.
\end{proof}

\begin{theorem*}[LADR 3.65]
  If $V$ and $W$ have the same finite dimension, then for $T\colon V\to W$ these are all the same condition:
  $T$ is injective; $T$ is surjective; $T$ is an isomorphism.
\end{theorem*}
\begin{proof}
  By rank--nullity, the null space has dimension $0$ exactly when the range has dimension $\dim V$, which equals $\dim W$.
\end{proof}
This is usually the fastest way to prove something is an isomorphism: show that $Tv=0$ forces $v=0$.

It also proves that equations have solutions. To show $Tp=q$ can be solved for every $q$, show $T$ is injective between spaces of the same dimension. The notes use this for polynomial interpolation (Ex.~34) and for $\big((x^2+5x+7)p\big)''=q$ (Ex.~35).

\begin{proposition*}[HW4 P3]
  Let $V$ be finite-dimensional and $T\colon V\to V$.
  $T$ is an isomorphism exactly when it sends a basis to a basis. It then sends every basis to a basis.
\end{proposition*}
\begin{proof}
  An isomorphism is injective, so it keeps lists independent, and surjective, so the image of a basis spans.
  Conversely, if $Tv_1,\dots,Tv_n$ is a basis, every vector is an output, so $T$ is surjective and therefore an isomorphism by LADR 3.65.
\end{proof}

\begin{proposition*}[LADR 3D]
  If $S$ and $T$ are invertible, then so is $ST$, and $(ST)^{-1}=T^{-1}S^{-1}$ (the order flips).

  If $ST$ is invertible, then $S$ is surjective and $T$ is injective.
\end{proposition*}
\begin{proof}
  If $S,T$ are invertible: $(ST)(T^{-1}S^{-1})=S(TT^{-1})S^{-1}=SS^{-1}=I$ and $(T^{-1}S^{-1})(ST)=T^{-1}(S^{-1}S)T=I$, so $ST$ is invertible with inverse $T^{-1}S^{-1}$.

  $S$ is surjective: any $v$ equals $S\big(T(ST)^{-1}v\big)$, which is an output of $S$.

  $T$ is injective: if $Tv=0$, then $v=(ST)^{-1}S(Tv)=(ST)^{-1}S(0)=0$.
\end{proof}

\begin{corollary*}[Prop.~33]
  If $V$ is finite-dimensional and $S,T\colon V\to V$, then $ST=I$ implies $TS=I$.
  So in this case you only need to check one side.
\end{corollary*}
\begin{proof}
  $ST=I$ makes $T$ injective. By LADR 3.65, $T$ is invertible. Then $S=S(TT^{-1})=(ST)T^{-1}=T^{-1}$.
\end{proof}

\begin{proposition*}[HW4 P1]
  The $2\times2$ matrix $A=\begin{psmallmatrix}a&b\\c&d\end{psmallmatrix}$, over any field, is invertible exactly when $ad-bc\ne0$. Then
  \[ A^{-1}=\frac{1}{ad-bc}\begin{pmatrix}d&-b\\-c&a\end{pmatrix}. \]
\end{proposition*}
\begin{proof}
  Multiplying out in either order gives $A\begin{psmallmatrix}d&-b\\-c&a\end{psmallmatrix}=\begin{psmallmatrix}d&-b\\-c&a\end{psmallmatrix}A=(ad-bc)I$. If $ad-bc\ne0$, divide by it.

  If $ad-bc=0$: when $a$ or $b$ is nonzero, $A$ sends the nonzero vector $(b,-a)$ to $0$, so it is not injective. When $a=b=0$, every output has first entry $0$, so $A$ is not surjective.
\end{proof}

\begin{proposition*}[HW4 P4]
  A product has rank at most the rank of each factor: $\operatorname{rank}(BC)\le\min(\operatorname{rank}B,\operatorname{rank}C)$.
\end{proposition*}
\begin{proof}
  Every output of $BC$ is an output of $B$, so $\operatorname{rank}(BC)\le\operatorname{rank}B$.
  The outputs of $BC$ are $B$ applied to the outputs of $C$, and a linear map cannot increase the dimension of a subspace, so $\operatorname{rank}(BC)\le\operatorname{rank}C$.
\end{proof}

\begin{proposition*}[Prop.~26]
  Once bases are fixed, sending $T$ to its matrix is an isomorphism from $\LL(V,W)$ to $\FF^{m,n}$.
  So $\dim\LL(V,W)=(\dim V)(\dim W)$.
\end{proposition*}

\begin{proposition*}[HW4 P7]
  Let $U$ be a subspace of a finite-dimensional $V$, and let $S\colon U\to V$ be injective.
  Then there is an invertible $T\colon V\to V$ that agrees with $S$ on $U$.
\end{proposition*}
\begin{proof}
  Take a basis $u_1,\dots,u_m$ of $U$ and add $w_1,\dots,w_k$ to get a basis of $V$.
  Because $S$ is injective, $Su_1,\dots,Su_m$ is independent. Add $y_1,\dots,y_k$ to it to get a basis of $V$.
  Define $T$ by $Tu_i=Su_i$ and $Tw_j=y_j$. It sends a basis to a basis, so it is invertible.
\end{proof}

\traps Do not write $T^{-1}$ until you have shown it exists.

LADR 3.65 and Prop.~33 need finite dimensions. On $\FF^\infty$, shifting left, $L(a_1,a_2,\dots)=(a_2,a_3,\dots)$, and shifting right, $R(a_1,a_2,\dots)=(0,a_1,a_2,\dots)$, satisfy $LR=I$ but $RL\ne I$. Differentiation on all polynomials is surjective but not injective.

A matrix that is not square has no two-sided inverse, but it can have a one-sided one. $\begin{psmallmatrix}1&2\end{psmallmatrix}$ has infinitely many right inverses.

If $ST$ is injective, you only learn that $T$ is injective, not $S$.

\subsection{Products}
\begin{definition*}
  The \vocab{product} $V_1\times\dots\times V_m$ is the set of lists $(v_1,\dots,v_m)$ with each $v_i$ in $V_i$.
  You add lists entry by entry, and multiply a list by a scalar by multiplying each entry.
\end{definition*}
The spaces do not need to be related. An element of $\PP_2(\RR)\times\RR^3$ is a pair: a polynomial and a vector.
$\RR^m\times\RR^n$ is really the same as $\RR^{m+n}$: just write the two lists one after the other.
Each $V_i$ sits inside the product as the lists that are $0$ except in slot $i$. The product is the direct sum of these copies.

\begin{proposition*}[Prop.~36]
  The dimensions add: $\dim(V_1\times\dots\times V_m)=\dim V_1+\dots+\dim V_m$.
\end{proposition*}
\begin{proof}
  For two spaces: if $v_1,\dots,v_n$ is a basis of $V$ and $w_1,\dots,w_k$ is a basis of $W$, then
  $(v_1,0),\dots,(v_n,0),(0,w_1),\dots,(0,w_k)$ is a basis of $V\times W$.
  It spans because $(v,w)=(v,0)+(0,w)$. It is independent because you can check each slot separately. For $m$ spaces, $V_1\times\dots\times V_m$ is $(V_1\times\dots\times V_{m-1})\times V_m$ with the same operations, so induct on $m$.
\end{proof}

\begin{proposition*}[Prop.~37]
  For subspaces $U_1,\dots,U_m$ of $V$, the map that adds up the entries, $(u_1,\dots,u_m)\mapsto u_1+\dots+u_m$, is a linear map from the product onto the sum $U_1+\dots+U_m$.
  It is injective exactly when the sum is direct.
  So, for finite-dimensional $U_i$, $\dim(U_1+\dots+U_m)\le\dim U_1+\dots+\dim U_m$, with equality exactly when the sum is direct.
\end{proposition*}

\traps The dimension of a product is a sum. The dimension of $\LL(V,W)$ is a product.

\subsection{Affine subsets}
\begin{definition*}
  An \vocab{affine subset} is a subspace moved away from the origin: $v+U$ is the set of all $v+u$ with $u$ in $U$.
  We say $v+U$ is \vocab{parallel} to $U$.
\end{definition*}
Examples:
\begin{itemize}
  \ii The line $x+y=1$ in $\RR^2$ is $(1,0)+\spn((-1,1))$: the line through the origin with slope $-1$, moved to pass through $(1,0)$.
  \ii The antiderivatives of $x^2$ are $\tfrac13x^3+C$ for any constant $C$. That is $\tfrac13x^3$ plus the subspace of constants.
  \ii The solutions of $Tx=b$ are either none, or one solution $x_0$ plus the null space: $x_0+\nul T$.
\end{itemize}
A map of the form $x\mapsto Tx+b$ with $T$ linear is called affine. It is linear only when $b=0$, because it sends $0$ to $b$.

\begin{proposition*}[LADR 3.101]
  For a subspace $U$, these three statements say the same thing:
  \begin{enumerate}
    \ii $v-w$ is in $U$;
    \ii $v+U=w+U$;
    \ii $v+U$ and $w+U$ have at least one point in common.
  \end{enumerate}
  So two shifted copies of $U$ are either identical or do not touch.
\end{proposition*}
\begin{proof}
  If $v-w$ is in $U$, then any $v+u$ can be rewritten as $w+\big((v-w)+u\big)$, which is in $w+U$. The same argument works the other way.

  If $v+U=w+U$, then $v=v+0$ is a point of both sets. And if the two sets share a point $v+u_1=w+u_2$, then $v-w=u_2-u_1$ is in $U$.
\end{proof}

\begin{proposition*}[HW4 P6, Prop.~38]
  $v+U$ is a subspace exactly when $v$ is in $U$, and then $v+U$ is just $U$.
  Each affine subset is parallel to only one subspace.
\end{proposition*}

\traps The set $v+U$ determines $U$, but not $v$. Any point of the set can play the role of $v$.
The solutions of $Tx=b$ with $b\ne0$ do not form a subspace, because $0$ is not a solution.

\subsection{Quotient spaces}
\begin{definition*}
  Let $U$ be a subspace of $V$. The \vocab{quotient} $V/U$ is the set of all the shifted copies $v+U$.
  You add and scale them using any representative:
  \[ (v+U)+(w+U)=(v+w)+U,\qquad \lambda(v+U)=\lambda v+U. \]
  The zero of $V/U$ is $0+U$, which is $U$ itself.
\end{definition*}
The idea: $V/U$ is $V$ where you stop telling apart two vectors that differ by something in $U$.

Example: take $U$ to be the $y$-axis in $\RR^2$. Then $(1,2)+U$ is the vertical line $x=1$, and so is $(1,4)+U$. Every vertical line is determined by its $x$-value, so $\RR^2/U$ is just a copy of $\RR$.

\textbf{Why ``well defined'' comes up.}
The same element of $V/U$ has many names: $v+U$ and $v'+U$ are equal whenever $v-v'$ is in $U$.
A formula that uses $v$ must give the same answer for every name.
For addition: if $v_1-v_2$ and $w_1-w_2$ are in $U$, then so is $(v_1+w_1)-(v_2+w_2)$, so the answer does not depend on the names. Scalar multiplication works the same way.

\begin{proposition*}[LADR 3.104, 3.105]
  The map $\pi(v)=v+U$ from $V$ to $V/U$ is linear and surjective, and it sends exactly the vectors of $U$ to zero.
  If $V$ is finite-dimensional, rank--nullity then gives $\dim V/U=\dim V-\dim U$.

  To get a basis: take a basis $u_1,\dots,u_m$ of $U$ and add $v_1,\dots,v_k$ to make a basis of $V$. Then $v_1+U,\dots,v_k+U$ is a basis of $V/U$.
\end{proposition*}

\begin{proposition*}[Prop.~42, maps out of a quotient]
  Let $T\colon V\to W$ be linear, and suppose $T$ sends everything in $U$ to $0$.
  Then $\tilde T(v+U)=Tv$ is a well-defined linear map from $V/U$ to $W$.
\end{proposition*}
\begin{proof}
  If $v+U=v'+U$, then $v-v'$ is in $U$, so $Tv-Tv'=T(v-v')=0$. Both names give the same output.
  Linear: $\tilde T\big((v+U)+(w+U)\big)=\tilde T\big((v+w)+U\big)=T(v+w)=Tv+Tw$, and $\tilde T\big(\lambda(v+U)\big)=T(\lambda v)=\lambda Tv$.
\end{proof}

\begin{theorem*}[First isomorphism theorem, LADR 3.107]
  For any linear $T\colon V\to W$, the map $v+\nul T\mapsto Tv$ is an isomorphism from $V/\nul T$ to $\range T$.
\end{theorem*}
\begin{proof}
  It is well defined by the previous proposition.
  It is surjective because every element of the range is some $Tv$.
  It is injective because if $Tv=0$, then $v$ is in the null space, so $v+\nul T$ is the zero element.
\end{proof}
Example: differentiation on polynomials sends exactly the constants to $0$. So ``polynomials, where you ignore the constant term'' is isomorphic to ``polynomials'', by $p\mapsto p'$. The inverse is the indefinite integral.

The second and third isomorphism theorems are $U/(U\cap W)\cong(U+W)/W$ and, for $W\subseteq U$, $(V/W)/(U/W)\cong V/U$. Both are proved the same way: write down a surjective map and apply the first theorem.

\traps Any map you define on $V/U$ by a formula in $v$ needs a check that it is well defined.
For example, $(x,y)+U\mapsto x$ works when $U$ is the $y$-axis but not when $U$ is the $x$-axis (Problem 12(c)).

$v+U$ is the zero of $V/U$ when $v$ is in $U$. It does not require $v=0$.

\clearpage
\section{Techniques}

Each technique is a statement you can use without proof, a note on why it is true, and problems that are hard without it and short with it. Work the problems before reading the solutions in Section~\ref{sec:sol}. Time on the exam is $1.25$ minutes per point; every technique here turns a ten-minute argument into two lines.

\newcommand{\why}{\par\noindent\textbf{Why.}\ }
\newcommand{\probs}{\par\medskip\noindent\textbf{Problems.}}

\subsection{Dimension pigeonhole}
\begin{theorem*}[T1]
  If $\dim V=\dim W<\infty$, a linear map $V\to W$ is injective if and only if it is surjective.
  So to prove ``$Tp=q$ has exactly one solution for every $q$'', prove only that $Tp=0$ forces $p=0$.
\end{theorem*}
\why Rank--nullity: $\dim\nul T=0$ exactly when $\dim\range T=\dim V=\dim W$.
The power of this is that injectivity is usually a one-line degree or roots argument, while surjectivity would require constructing a solution.
\probs
\begin{enumerate}[(a)]
  \ii Prove that for every $q\in\PP_n(\RR)$ there is exactly one $p\in\PP_n(\RR)$ with $p+p'+p''=q$.
  \ii Prove that for every $q\in\PP_n(\RR)$ there is a $p\in\PP_n(\RR)$ with $\big((x^2+1)p\big)''=q$. (Notes Ex.~35, varied.) Then show that the similar-looking map $p\mapsto(x^2+1)p''$ is \emph{not} surjective on $\PP_n(\RR)$.
  \ii Let $a_0,\dots,a_n$ be distinct reals. Prove that $p\mapsto(p(a_0),\dots,p(a_n))$ is an isomorphism $\PP_n(\RR)\to\RR^{n+1}$, and deduce that any function on $\{a_0,\dots,a_n\}$ is given by a polynomial of degree $\le n$.
\end{enumerate}

\subsection{Forced relations}
\begin{theorem*}[T2]
  Any $n+1$ vectors in an $n$-dimensional space are dependent.
  If subspaces $U_1,\dots,U_k$ of $V$ have $\dim U_1+\dots+\dim U_k>(k-1)\dim V$, then $U_1\cap\dots\cap U_k\ne\{0\}$.
\end{theorem*}
\why The first is the replacement lemma. The second is $\dim(U\cap W)\ge\dim U+\dim W-\dim V$ applied $k-1$ times.
Use it whenever you need a nonzero object satisfying constraints: count dimensions, and existence is free.
The pattern: the objects you want form a subspace (or the constraints cut out a subspace); if its dimension is forced to be positive, a nonzero element exists, and you never construct it. Problem (a) below is the model: you get a polynomial relation for $T$ with no idea what it is, purely because $n^2+1>n^2$.
\probs
\begin{enumerate}[(a)]
  \ii Let $\dim V=n$ and $T\in\LL(V)$. Prove that there is a nonzero polynomial $p$ of degree at most $n^2$ with $p(T)=0$.
  \ii Let $T\colon V\to W$ with $\dim V>\dim W$. Prove that some nonzero $v$ has $Tv=0$, without writing the words ``rank--nullity''.
  \ii Also: Problem 4(a) (three $7$-dimensional subspaces of a $10$-dimensional space).
\end{enumerate}

\subsection{Restrict and count}
\begin{theorem*}[T3]
  For any linear $T\colon V\to W$ and any finite-dimensional subspace $X\subseteq V$,
  \[ \dim X=\dim(X\cap\nul T)+\dim T(X). \]
\end{theorem*}
\why This is rank--nullity for the map $T|_X\colon X\to W$. Its null space is the set of $x\in X$ with $Tx=0$, which is $X\cap\nul T$; its range is $T(X)$.

\medskip\noindent\textbf{What it is for.} Every formula in this course of the shape
\[ \dim(\text{something})=\dim(\text{a null space})+\dim(\text{a range}) \]
is this identity. The skill is choosing $X$ and $T$. The recipe:
\begin{enumerate}
  \ii The left side names $X$: it is the space whose dimension you want.
  \ii Ask what natural map is defined on $X$. It is almost always a map you already have, restricted to $X$.
  \ii Compute $X\cap\nul T$ and $T(X)$ by unpacking definitions. Rank--nullity then \emph{is} the formula.
\end{enumerate}
On the practice midterm, Hammond handed you $X=T^{-1}(U)$ and the map $S=T|_X$ and walked you through steps 2 and 3. On the real exam he may hand you a different $X$, or none. Here are the other three cases worked in full, so you can see the recipe run without help. These three derivations are also one of the two abstract questions on each of Practice Midterms A, B and C (A Q4, B Q2, C Q2); if you want to take those exams cold, skip to the summary table now and come back afterwards.

\medskip\noindent\textbf{Derivation 1: $\dim\nul(ST)$.} Want $\dim\nul(ST)$, so $X=\nul(ST)\subseteq U$ for $T\colon U\to V$, $S\colon V\to W$. The natural map on $X$ is $T$.
Null space of $T|_X$: vectors in $X$ with $Tu=0$. Every $u$ with $Tu=0$ has $STu=0$, so $u\in X$ automatically: $X\cap\nul T=\nul T$.
Range of $T|_X$: vectors $Tu$ with $S(Tu)=0$, i.e.\ vectors that are in $\range T$ and in $\nul S$: $T(X)=\nul S\cap\range T$. (Both inclusions: if $y=Tu$ with $Sy=0$ then $y$ is in both; if $y\in\nul S\cap\range T$, write $y=Tu$, then $STu=0$ so $u\in X$.)
Therefore
\[ \dim\nul(ST)=\dim\nul T+\dim(\nul S\cap\range T)\le\dim\nul T+\dim\nul S. \]

\medskip\noindent\textbf{Derivation 2: $\dim(U+W)$.} Want $\dim(U+W)$. There is no map on $U+W$ itself, so instead take a space that maps \emph{onto} $U+W$: $X=U\times W$ with $\Gamma(u,w)=u+w$. Then $\dim X=\dim U+\dim W$ (product), $\range\Gamma=U+W$, and $\nul\Gamma=\{(u,-u):u\in U\cap W\}$, which is a copy of $U\cap W$. Rank--nullity:
\[ \dim U+\dim W=\dim(U\cap W)+\dim(U+W). \]
This is SLC problem 9, and it also shows why $U\times W\cong U\oplus W$ when the intersection is zero: then $\Gamma$ is injective.

\medskip\noindent\textbf{Derivation 3: rank of a product.} Want $\rank(ST)=\dim\range(ST)$. Now $X=\range T$ and the map is $S$: $S(\range T)=\range(ST)$, and $\range T\cap\nul S$ is the null space. So
\[ \rank T=\dim(\nul S\cap\range T)+\rank(ST). \]
Read it two ways. Since the middle term is $\ge0$: $\rank(ST)\le\rank T$. Since the middle term is $\le\dim\nul S=\dim V-\rank S$: $\rank(ST)\ge\rank S+\rank T-\dim V$ (Sylvester). Equality $\rank(ST)=\rank T$ holds exactly when $S$ is injective on $\range T$.

\medskip\noindent\textbf{Summary table.}
\begin{center}
\renewcommand{\arraystretch}{1.2}
\small\begin{tabular}{lll}
$X$ & $T(X)$ & Formula \\ \hline
$T^{-1}(U)$ & $U\cap\range T$ & $\dim T^{-1}(U)=\dim\nul T+\dim(U\cap\range T)$ \\
$\nul(ST)$ & $\nul S\cap\range T$ & $\dim\nul(ST)=\dim\nul T+\dim(\nul S\cap\range T)$ \\
$\range T$, map $S$ & $\range(ST)$ & $\rank T=\dim(\nul S\cap\range T)+\rank(ST)$ \\
$U\times W$, map $\Gamma$ & $U+W$ & $\dim U+\dim W=\dim(U\cap W)+\dim(U+W)$
\end{tabular}
\end{center}
\probs
\begin{enumerate}[(a)]
  \ii Prove $\dim\nul(T^2)\le2\dim\nul T$ and $\rank(T^2)\ge2\rank T-\dim V$ for $T\in\LL(V)$.
  \ii Let $T\colon V\to W$ and let $X\subseteq V$ be a subspace with $X\cap\nul T=\{0\}$. Prove that $T$ restricted to $X$ is an isomorphism onto $T(X)$, and that $\dim X\le\rank T$.
  \ii For $T\in\LL(V)$ prove $\dim\nul(T^2)=\dim\nul T+\dim(\nul T\cap\range T)$, and conclude that $\nul(T^2)=\nul T$ if and only if $\nul T\cap\range T=\{0\}$.
  \ii Let $U\subseteq V$ and $T\in\LL(V,W)$. Prove $\dim T(U)\ge\dim U-\dim\nul T$, and find $T$, $U$ with equality and with strict inequality.
\end{enumerate}

\subsection{Roots versus degree}
\begin{theorem*}[T4]
  A polynomial of degree $\le n$ with $n+1$ distinct roots is zero.
  If $p(a)=p'(a)=\dots=p^{(m-1)}(a)=0$, then $(x-a)^m$ divides $p$ (expand $p$ in powers of $x-a$; the first $m$ coefficients are $p^{(j)}(a)/j!$ and vanish), so a root of multiplicity $m$ counts $m$ times.
  Consequently, over $\RR$, a map $\PP_n(\RR)\to\RR^{n+1}$ that records at each of several points $a$ the values $p(a),p'(a),\dots,p^{(m_a-1)}(a)$, all derivatives up to some order, with $\sum_a m_a=n+1$, is injective, hence an isomorphism by T1. Skipping an order at a point breaks this: see problem (b).
\end{theorem*}
\why Each root $a$ contributes a factor $(x-a)$; a nonzero polynomial cannot have more linear factors than its degree.
\probs
\begin{enumerate}[(a)]
  \ii Prove that $p\mapsto(p(0),p'(0),p(1),p'(1))$ is an isomorphism $\PP_3(\RR)\to\RR^4$.
  \ii For which $a\in\RR$ is $p\mapsto(p(0),p(1),p'(a))$ an isomorphism $\PP_2(\RR)\to\RR^3$? This is Problem 8(a); here, say in one sentence which hypothesis of T4 fails at the bad $a$.
  \ii Let $a_0,\dots,a_n$ be distinct and $L_i(x)=\prod_{j\ne i}\frac{x-a_j}{a_i-a_j}$. Prove $L_0,\dots,L_n$ is a basis of $\PP_n(\RR)$ by evaluating a relation at each $a_k$, and write the interpolating polynomial through $(a_i,b_i)$ in this basis.
\end{enumerate}

\subsection{Degree filtration}
\begin{theorem*}[T5]
  Let $T\colon\PP_n\to\PP_n$ not raise degree. In the monomial basis, $\MM(T)$ is upper triangular; the diagonal entry in column $k$ is the coefficient of $x^k$ in $T(x^k)$.
  $T$ is invertible if and only if every diagonal entry is nonzero.
  If $T$ strictly lowers degree, then $T^{n+1}=0$.
  Nonzero polynomials of distinct degrees are independent.
\end{theorem*}
\why $T(x^k)$ only involves $1,\dots,x^k$. A triangular matrix with nonzero diagonal has independent columns: in a relation among the columns, the last column with a nonzero coefficient is the only one with an entry in its diagonal row, so that coefficient is $0$; repeat, exactly as in Ex.~18. If some diagonal entry is $0$, say in column $k$, then columns $1,\dots,k$ all lie in the span of the first $k-1$ coordinate vectors, so they are dependent and $T$ is not injective. If $T$ strictly lowers degree, $T^{n+1}$ kills every monomial.

The underlying principle: once bases are fixed, $T$ is injective iff the columns of $\MM(T)$ are independent, $\rank T=\rank\MM(T)$, and $T$ is invertible iff $\MM(T)$ is. So a basis in which $\MM(T)$ is triangular or diagonal lets you read everything off the diagonal. The monomial basis does this for any map that does not raise degree, including the shift $p(x)\mapsto p(x+1)$; Midterm B Q3 computes the shift in the monomial basis and in $1,x,x(x-1)$, and the diagonal is all $1$'s either way.
\probs
\begin{enumerate}[(a)]
  \ii Prove $p\mapsto p(2x)$ and $p\mapsto p(x+1)$ are isomorphisms of $\PP_n(\RR)$ by reading the diagonal.
  \ii Let $\Delta p=p(x+1)-p(x)$ on $\PP_n(\RR)$. Prove $\nul\Delta$ is the constants and $\range\Delta=\PP_{n-1}(\RR)$, using only degrees.
  \ii Let $Tp=xp'-p$ on $\PP_n(\RR)$. Find $\nul T$ and $\range T$ in one line each.
\end{enumerate}

\subsection{Nilpotent inverse}
\begin{theorem*}[T6]
  If $N^k=0$, then $I-N$ is invertible with $(I-N)^{-1}=I+N+N^2+\dots+N^{k-1}$, and $(I+N)^{-1}=I-N+N^2-\dots\pm N^{k-1}$.
  This holds on infinite-dimensional spaces too, as long as each vector is killed by some power of $N$.
\end{theorem*}
\why $(I-N)(I+N+\dots+N^{k-1})=I-N^k=I$, and the factors commute.
Differentiation on $\PP_n$ is the standard $N$: it turns $p\pm p'\pm p''=q$ into a formula.
\probs
\begin{enumerate}[(a)]
  \ii Solve $p-p'+p''-p'''=x^3$ in $\PP_3(\RR)$ without solving a linear system.
  \ii Let $N$ be the operator $p\mapsto xp'$ minus $p\mapsto p$, i.e.\ $Np=xp'-p$, on $\PP_2(\RR)$. Is $I-N$ invertible? Decide using T5, and explain why T6 does not apply.
  \ii Also: Problem 7(d), 7(e); Midterm B Q4(c),(d).
\end{enumerate}

\subsection{Polynomial identities in \texorpdfstring{$T$}{T}}
\begin{theorem*}[T7]
  Suppose $a_0I+a_1T+\dots+a_kT^k=0$ with $a_0\ne0$. Then $T$ is invertible and
  $T^{-1}=-\frac1{a_0}\big(a_1I+a_2T+\dots+a_kT^{k-1}\big)$.
  If instead $p(T)=0$ and $p(x)=(x-\lambda)q(x)$, then on a finite-dimensional space either $T-\lambda I$ is not injective or $q(T)=0$.
\end{theorem*}
\why Move $a_0I$ to one side and factor out $T$. For the second part: $(T-\lambda I)q(T)=0$; if $T-\lambda I$ were injective, it would be invertible, and multiplying by its inverse gives $q(T)=0$.
\probs
\begin{enumerate}[(a)]
  \ii Also: Problem 10(d) ($T^2-3T+2I=0$ over $\RR$).
  \ii $T^3=I$. Prove $T$ is invertible and find $T^{-1}$. If also $T\ne I$, prove that $T^2+T+I$ is not injective.
  \ii Over $\FF_2$, $T^2+T=0$ tells you nothing about invertibility. Give a $T$ satisfying it that is invertible and one that is not.
\end{enumerate}

\subsection{\texorpdfstring{$AB=0$}{AB=0}}
\begin{theorem*}[T8]
  If $AB=0$, then $\range B\subseteq\nul A$. Hence: if $B\ne0$ then $A$ is not injective; if $A\ne0$ then $B$ is not surjective; and for operators on $V$, $\rank A+\rank B\le\dim V$.
\end{theorem*}
\why $A(Bv)=0$ for every $v$. Then $\rank B\le\dim\nul A=\dim V-\rank A$.
\probs
\begin{enumerate}[(a)]
  \ii Can $S,T\in\LL(\RR^4)$ both have rank $2$ with $ST=0$? Give one or prove not. (The $\RR^3$ version is Problem 5(d).)
  \ii $T^2=0$ on $\RR^4$. Give a $T$ with $\rank T=2$, and prove no $T$ with $T^2=0$ has larger rank. (General bound: Problem 6(a), Midterm B Q4.)
  \ii $P^2=P$ and $P\ne I$. Prove $P$ is not injective. $P^2=P$ and $P\ne0$. Prove $I-P$ is not injective.
\end{enumerate}

\subsection{Define it on a basis}
\begin{theorem*}[T9]
  To build a linear map with prescribed behaviour: choose a basis of the domain, adapted to the problem (a basis of a given subspace, extended), assign each basis vector freely, and cite LADR 3.4.
  To handle vectors outside a subspace, extend a basis of the subspace and send the new vectors wherever the property needs, usually to $0$ or to an independent list.
\end{theorem*}
\why LADR 3.4: the assignment on a basis extends uniquely to a linear map. Every ``there exists a linear map'' problem in this course is this move.
\probs
\begin{enumerate}[(a)]
  \ii Let $v_1,\dots,v_k$ be independent in $V$ and $w_1,\dots,w_k$ arbitrary in $W$. Prove there is a linear $T$ with $Tv_i=w_i$. Show by example that independence cannot be dropped.
  \ii Let $U\subseteq V$ be a subspace. Prove there is $T\in\LL(V)$ with $\range T=U$, and one with $\nul T=U$, and one with both ($\nul T=\range T=U$) if and only if $\dim U=\frac12\dim V$.
  \ii Let $T\in\LL(V)$ and suppose $Tv\in\spn(v)$ for every $v\in V$. Prove $T=\lambda I$ for a single scalar $\lambda$. (This is the second half of Problem 10(c); the first half builds $S$ to force the hypothesis.)
  \ii Also: Problem 6(c) (prescribed null space into a smaller space), Problem 11(e), Hammond's practice midterm Q2(b) and Practice Midterm A Q2(b).
\end{enumerate}

\subsection{Idempotents and involutions split the space}
\begin{theorem*}[T10]
  If $P^2=P$, then $V=\nul P\oplus\range P$, with $v=(v-Pv)+Pv$.
  If $J^2=I$ and $2\ne0$ in $\FF$, then $V=\{Jv=v\}\oplus\{Jv=-v\}$, with $v=\tfrac12(v+Jv)+\tfrac12(v-Jv)$.
\end{theorem*}
\why $P(v-Pv)=Pv-P^2v=0$, and the intersection is zero: if $v=Pw$ and $Pv=0$ then $v=Pw=P^2w=Pv=0$. For $J$, apply this to $P=\tfrac12(I+J)$.
The involutions you will meet: transpose ($A\mapsto A^T$), reflection ($f(x)\mapsto f(-x)$), and swapping coordinates.
\probs
\begin{enumerate}[(a)]
  \ii Prove $\RR^{n,n}=\{\text{symmetric}\}\oplus\{\text{skew}\}$ and give both dimensions. (The $2\times2$ case is Midterm B Q1(a).)
  \ii Prove $\PP_n(\RR)=\{\text{even}\}\oplus\{\text{odd}\}$, and find both dimensions.
  \ii $P^2=P$. Prove $I-P$ is also idempotent, and that $\nul(I-P)=\range P$ and $\range(I-P)=\nul P$.
  \ii Give $T\in\LL(\RR^2)$ with $\RR^2\ne\nul T\oplus\range T$, and explain which hypothesis of T10 it violates.
\end{enumerate}

\subsection{Two names for the same thing}
\begin{theorem*}[T11]
  Direct sums: to prove uniqueness of $v=u+w$, subtract two representations; the difference $u-u'=w'-w$ lies in $U\cap W$.
  Quotients: a formula $f(v+U)$ is well defined if and only if $f(u)=0$ for every $u\in U$ (for linear $f$). Check $f$ on $U$, not on pairs of representatives.
\end{theorem*}
\why If $f$ is linear and kills $U$, then $v-v'\in U$ gives $f(v)=f(v')$. Conversely $0+U=u+U$ forces $f(u)=f(0)=0$.
\probs
\begin{enumerate}[(a)]
  \ii On $\RR^2/\spn((1,1))$, which of these are well defined: $(x,y)+U\mapsto x$; $\mapsto x-y$; $\mapsto x+y$; $\mapsto x^2-y^2$? One line each.
  \ii Let $U\subseteq V$ and let $f\colon V\to\FF$ be linear with $f(U)=0$. Prove that $f$ induces a well-defined $\bar f\colon V/U\to\FF$ and that $\bar f$ is injective if and only if $U=\nul f$. (General version: Problem 12(d).)
  \ii $U\oplus W=V$. Define $P\colon V\to V$ by $P(u+w)=u$. Prove $P$ is well defined and linear, $P^2=P$, $\range P=U$, $\nul P=W$.
\end{enumerate}

\subsection{Counting over \texorpdfstring{$\FF_p$}{F\_p}}
\begin{theorem*}[T12]
  A $d$-dimensional space over $\FF_p$ has $p^d$ elements.
  A consistent system $Ax=b$ has exactly $p^{\,n-\rank A}$ solutions ($n$ unknowns).
  $|GL_n(\FF_p)|=(p^n-1)(p^n-p)\cdots(p^n-p^{n-1})$: choose the columns one at a time, each outside the span of the previous ones.
  The number of ordered bases of $\FF_p^n$ is the same product.
\end{theorem*}
\why Solutions form $x_0+\nul A$, a translate of a $(n-\rank A)$-dimensional space. A matrix is invertible iff its columns are a basis, and column $k+1$ may be any vector outside a $k$-dimensional span, which has $p^k$ elements.
\probs
\begin{enumerate}[(a)]
  \ii Count the invertible $3\times3$ matrices over $\FF_2$.
  \ii Count the $2$-dimensional subspaces of $\FF_2^3$. (Count ordered pairs of independent vectors in the whole space and ordered bases of one subspace, then divide.)
  \ii Over $\FF_3$, how many solutions does $x+y+z=1,\ x-y=2$ have? List them.
  \ii How many polynomials in $\PP_4(\FF_2)$ define the zero function on $\FF_2$? (Use $x^2=x$ as functions.)
\end{enumerate}

\subsection{Universal counterexamples}
\begin{theorem*}[T13]
  Six objects refute almost every false statement in this course. Know what each one kills.
\end{theorem*}
\begin{center}
\renewcommand{\arraystretch}{1.25}
\begin{tabular}{p{0.34\textwidth}p{0.6\textwidth}}
\hline
$N(x,y)=(y,0)$ on $\RR^2$ & $\nul N=\range N$; $N^2=0$, $N\ne0$; $V\ne\nul N\oplus\range N$; rank--nullity gives no decomposition. \\
Shifts $L,R$ on $\FF^\infty$ & $LR=I\ne RL$; $R$ injective not surjective; $L$ surjective not injective; T1 needs finite dimension. \\
$D$ on $\PP(\RR)$ & surjective, not injective; $\int$ is a right inverse but not a left inverse. \\
Three lines in $\RR^2$ & pairwise intersections zero but the sum is not direct; no inclusion--exclusion for three subspaces. \\
$\FF_2$ & $2=0$ kills every ``divide by $2$''; $(1,1,0),(0,1,1),(1,0,1)$ dependent; $x^2=x$ as functions; $T^2+T=0$ says nothing about invertibility. \\
$E_{11}$, $E_{22}$ & non-invertible, but their sum $I$ is invertible: singular matrices are not a subspace. \\
\hline
\end{tabular}
\end{center}
\probs
\begin{enumerate}[(a)]
  \ii For each statement, decide true or false and name the counterexample from the table if false: (i) $\nul T\cap\range T=\{0\}$ for every operator; (ii) $ST=I$ implies $TS=I$ for operators on any space; (iii) an injective operator on a finite-dimensional space is surjective; (iv) if $U_i\cap U_j=\{0\}$ for all $i\ne j$ then $U_1+U_2+U_3$ is direct.
  \ii Find an operator $T$ on $\RR^4$ with $T^3=0$ but $T^2\ne0$. What are $\rank T$, $\rank T^2$?
\end{enumerate}

\subsection{Facts you may quote}
Cite by name; do not reprove.
Injective $\iff\nul T=\{0\}$.
Rank--nullity.
Replacement lemma: independent length $\le$ spanning length; so bases have equal length, $\dim U\le\dim V$, and independent lists extend to bases.
$n$ independent vectors in an $n$-dimensional space are a basis; a subspace of full dimension is the whole space.
Differentiation $\PP_n\to\PP_{n-1}$ is surjective with null space the constants.
A linear map is determined by a basis (LADR 3.4).
Injective maps preserve independence; surjective maps preserve spanning.
$U+W$ direct $\iff U\cap W=\{0\}$; $\dim(U+W)=\dim U+\dim W-\dim(U\cap W)$.
$ST=I\Rightarrow TS=I$ for operators on a finite-dimensional space.
$\dim(V\times W)=\dim V+\dim W$; $\dim\LL(V,W)=(\dim V)(\dim W)$; $\dim V/U=\dim V-\dim U$.

\section{Practice midterms}

Three exams in the format of the practice midterm: 4 questions, 40 points, 50 minutes, points per part in parentheses. Two questions are computations with a short proof attached; two are abstract proofs where the parts lay out the route and, usually, the last part applies the result to a concrete example, often polynomials. Exams A and B each have a question using a nonstandard basis. You may quote any result from the notes if you state it.
A is at the level of the practice exam. B and C are harder, at the level of SLC problem 9. Do each one with a timer before reading the solutions in Section~\ref{sec:sol}.
Notation as on the real exam: $\PP_n(\RR)$ is polynomials of degree at most $n$; all bases are ordered.

\clearpage\begingroup\subsection{Practice Midterm A}

\medskip\noindent\textbf{1. A direct sum in $\PP_2(\RR)$.} Let $U=\{p\in\PP_2(\RR):p(0)=0\}$ and $W=\spn(1+x)$.
\begin{enumerate}[a)]
  \ii (6) Prove that $\PP_2(\RR)=U\oplus W$.
  \ii (2) Write $2+3x+x^2$ in the form $u+w$ with $u\in U$ and $w\in W$.
\end{enumerate}

\medskip\noindent\textbf{2. A right inverse.} Let $V$ and $W$ be finite-dimensional and let $T\colon V\to W$ be a surjective linear map.
\begin{enumerate}[a)]
  \ii (4) Choose a basis $w_1,\dots,w_m$ of $W$ and choose $v_1,\dots,v_m\in V$ with $Tv_j=w_j$. Prove that $v_1,\dots,v_m$ is linearly independent in $V$.
  \ii (4) Prove that there exists a linear map $S\colon W\to V$ such that $TS=I_W$.
  \ii (2) Give an example with $V=\RR^2$, $W=\RR$ in which two different maps $S$ satisfy $TS=I_W$.
\end{enumerate}

\medskip\noindent\textbf{3. A matrix of a linear map.} Define $T\colon\PP_2(\RR)\to\RR^3$ by
\[ T(p)=\big(p(0),\,p(1),\,p(-1)\big). \]
Use the basis $B=(1,x,x^2)$ of the domain and the basis $C=\big((1,0,0),(1,1,0),(1,1,1)\big)$ of the codomain.
\begin{enumerate}[a)]
  \ii (6) Find $\MM(T)$ with respect to $B$ and $C$.
  \ii (2) Prove that $T$ is an isomorphism.
  \ii (2) Find the polynomial $p$ with $T(p)=(0,1,0)$.
\end{enumerate}

\medskip\noindent\textbf{4. The null space of a composition.} Let $U,V$ be finite-dimensional, and let $T\colon U\to V$ and $S\colon V\to W$ be linear. Consider the map
\[ R\colon\nul(ST)\longrightarrow\nul S\cap\range T,\qquad R(u)=Tu. \]
\begin{enumerate}[a)]
  \ii (3) Prove that $R(u)\in\nul S\cap\range T$ for every $u\in\nul(ST)$. Then prove that $R$ is linear and surjective.
  \ii (2) Prove that $\nul R=\nul T$.
  \ii (3) Apply the fundamental theorem of linear maps to $R$ to prove that
  \[ \dim\nul(ST)=\dim\nul T+\dim\big(\nul S\cap\range T\big). \]
  \ii (4) Let $T\colon\PP_3(\RR)\to\PP_3(\RR)$ be differentiation and let $S\colon\PP_3(\RR)\to\RR^2$ be $S(p)=(p(0),p(1))$. Use part (c) to determine $\dim\nul(ST)$, and then give a basis of $\nul(ST)$.
\end{enumerate}

\endgroup
\clearpage\begingroup\subsection{Practice Midterm B}

\medskip\noindent\textbf{1. Symmetric and skew matrices.} In $V=\RR^{2,2}$ let
\[ U=\{A:A^T=A\},\qquad W=\{A:A^T=-A\}. \]
\begin{enumerate}[a)]
  \ii (4) Prove that $V=U\oplus W$, and state $\dim U$ and $\dim W$.
  \ii (2) Write $\begin{psmallmatrix}1&2\\3&4\end{psmallmatrix}$ in the form $u+w$ with $u\in U$, $w\in W$.
  \ii (2) Define $T\colon V\to V$ by $T(A)=A-A^T$. Determine $\nul T$ and $\range T$.
\end{enumerate}

\medskip\noindent\textbf{2. Sum and product.} Let $U,W$ be subspaces of a finite-dimensional space $V$. Define
\[ \Gamma\colon U\times W\longrightarrow V,\qquad \Gamma(u,w)=u+w. \]
\begin{enumerate}[a)]
  \ii (3) Prove that $\Gamma$ is linear and that $\range\Gamma=U+W$.
  \ii (3) Prove that $\nul\Gamma=\{(u,-u):u\in U\cap W\}$, and that $u\mapsto(u,-u)$ is an isomorphism $U\cap W\to\nul\Gamma$.
  \ii (2) Deduce that $\dim(U+W)=\dim U+\dim W-\dim(U\cap W)$. You may use $\dim(U\times W)=\dim U+\dim W$.
  \ii (2) Prove that if $U\cap W=\{0\}$, then $U\times W$ is isomorphic to $U\oplus W$.
\end{enumerate}

\medskip\noindent\textbf{3. A shift of the variable.} Define $T\colon\PP_2(\RR)\to\PP_2(\RR)$ by $(Tp)(x)=p(x+1)$.
\begin{enumerate}[a)]
  \ii (4) Find $\MM(T)$ with respect to the standard basis $(1,x,x^2)$ on both sides.
  \ii (3) Find $\MM(T)$ with respect to the basis $B=\big(1,\,x,\,x(x-1)\big)$ on both sides.
  \ii (3) Prove that $T$ is an isomorphism, and give a formula for $T^{-1}$.
\end{enumerate}

\medskip\noindent\textbf{4. An operator with $T^2=0$.} Let $V$ be finite-dimensional and let $T\in\LL(V)$ satisfy $T^2=0$.
\begin{enumerate}[a)]
  \ii (3) Prove that $\range T\subseteq\nul T$.
  \ii (3) Prove that $\dim\range T\le\frac12\dim V$.
  \ii (3) Prove that $I-T$ is invertible, and find its inverse.
  \ii (3) Let $T\colon\PP_3(\RR)\to\PP_3(\RR)$ be $Tp=p''$. Verify $T^2=0$, find $\range T$ and $\nul T$, and solve $p-p''=x^3$ for $p\in\PP_3(\RR)$.
\end{enumerate}

\endgroup
\clearpage\begingroup\subsection{Practice Midterm C}

\medskip\noindent\textbf{1. Polynomials vanishing at a point.} Let $U=\{p\in\PP_3(\RR):p(2)=0\}$.
\begin{enumerate}[a)]
  \ii (3) Prove that $U$ is a subspace and determine $\dim U$.
  \ii (3) Give a basis of $U$, with proof.
  \ii (2) Give a subspace $W$ with $\PP_3(\RR)=U\oplus W$.
\end{enumerate}

\medskip\noindent\textbf{2. Rank of a composition.} Let $U,V,W$ be finite-dimensional and let $T\colon U\to V$, $S\colon V\to W$ be linear.
\begin{enumerate}[a)]
  \ii (3) Prove that $\range(ST)=S(\range T)$, and that $\range(ST)\subseteq\range S$.
  \ii (3) Prove that $\dim\range(ST)\le\min\{\dim\range S,\dim\range T\}$.
  \ii (4) Prove that $\dim\range(ST)\ge\dim\range S+\dim\range T-\dim V$.
\end{enumerate}

\medskip\noindent\textbf{3. An interpolation map.} Define $T\colon\PP_2(\RR)\to\RR^3$ by $T(p)=\big(p(0),\,p(1),\,p'(1)\big)$.
\begin{enumerate}[a)]
  \ii (4) Find $\MM(T)$ with respect to the standard bases.
  \ii (3) Prove that $T$ is an isomorphism.
  \ii (3) Find the unique $p\in\PP_2(\RR)$ with $p(0)=1$, $p(1)=0$, $p'(1)=2$.
\end{enumerate}

\medskip\noindent\textbf{4. A quotient by the null space.} Let $V,W$ be finite-dimensional and $T\colon V\to W$ linear. Define
\[ \tilde T\colon V/\nul T\longrightarrow\range T,\qquad \tilde T(v+\nul T)=Tv. \]
\begin{enumerate}[a)]
  \ii (3) Prove that $\tilde T$ is well defined: if $v+\nul T=v'+\nul T$, then $Tv=Tv'$.
  \ii (3) Prove that $\tilde T$ is linear, injective and surjective.
  \ii (2) Deduce that $\dim(V/\nul T)=\dim\range T$.
  \ii (4) Let $T\colon\PP_2(\RR)\to\RR$ be $T(p)=p(1)-p(0)$. Find a basis of $\nul T$, a basis of $\PP_2(\RR)/\nul T$, and decide whether $x+\nul T=x^2+\nul T$.
\end{enumerate}

\endgroup
\clearpage
\section{Problem bank}
Harder, multi-part problems in homework style. Starred problems are the most likely exam material. Solutions are in Section~\ref{sec:sol}.

\begin{problem}[$\star$ Images and preimages]
  Let $f\colon X\to Y$ be a function.
  \begin{enumerate}[(a)]
    \ii Prove $f$ is injective if and only if $f^{-1}(f(U))=U$ for every $U\subseteq X$.
    \ii Prove $f$ is surjective if and only if $f(f^{-1}(V))=V$ for every $V\subseteq Y$.
    \ii Prove $f(A\cap B)\subseteq f(A)\cap f(B)$, with equality for all $A,B$ if and only if $f$ is injective.
    \ii Let $T\colon V\to W$ be linear, $U\subseteq V$ and $X\subseteq W$ subspaces. Prove
    \[ T^{-1}(T(U))=U+\nul T, \qquad T(T^{-1}(X))=X\cap\range T. \]
    \ii For $V$ finite-dimensional, prove $\dim T(U)=\dim U-\dim(U\cap\nul T)$.
    Then prove $T^{-1}(T(U))=U$ for every subspace $U$ if and only if $T$ is injective.
  \end{enumerate}
\end{problem}

\begin{problem}[Finite fields and polynomial functions]
  Let $p$ be prime. Here $\PP_n(\Fp)$ means formal polynomials of degree $\le n$ (basis $1,x,\dots,x^n$), and $\SF$ is the space of all functions $\Fp\to\Fp$.
  \begin{enumerate}[(a)]
    \ii Prove $\ZZ/6\ZZ$ is not a field.
    \ii Let $\FF$ be a field with $q$ elements and $a\ne0$.
    Prove $x\mapsto ax$ is a bijection $\FF\to\FF$, and deduce $a^{q-1}=1$.
    Conclude $c^p=c$ for all $c\in\Fp$.
    \ii Find $\dim\SF$.
    Let $E\colon\PP_n(\Fp)\to\SF$ send a polynomial to the function it defines.
    Determine exactly for which $n$ the map $E$ is injective, surjective, an isomorphism.
    For $n=p$, give a nonzero element of $\nul E$.
    \ii Over $\FF_2$, expand $(x+1)^4$.
    Which polynomial of degree $\le 1$ defines the same function?
    Which hypothesis of the interpolation argument (Ex.~34) fails when $|\FF|<n+1$?
  \end{enumerate}
\end{problem}

\begin{problem}[Independence depends on the field]
  Let $v_1,\dots,v_n$ be linearly independent, $n\ge2$.
  \begin{enumerate}[(a)]
    \ii Prove or disprove: $v_1-v_2,\ v_2-v_3,\ \dots,\ v_{n-1}-v_n,\ v_n-v_1$ is linearly independent.
    \ii Let $w_j=v_j+v_{j+1}$ for $j<n$ and $w_n=v_n+v_1$.
    Determine exactly, in terms of $n$ and $\FF$, when $w_1,\dots,w_n$ is independent.
    In that case, write $v_1$ in terms of the $w_j$.
    \ii Suppose $v_1+w,\dots,v_n+w$ is linearly dependent. Prove $w\in\spn(v_1,\dots,v_n)$.
    \ii Prove or disprove: if $\spn(u_1,\dots,u_n)=\spn(v_1,\dots,v_n)$, then $u_1,\dots,u_n$ is independent.
  \end{enumerate}
\end{problem}

\begin{problem}[Dimension counting]
  \leavevmode
  \begin{enumerate}[(a)]
    \ii Let $\dim V=10$ and $U,W,X$ subspaces of dimension $7$.
    Prove $U\cap W\cap X\ne\{0\}$, and show $\dim(U\cap W\cap X)=1$ is possible.
    \ii If $U_1,\dots,U_k$ are subspaces of an $n$-dimensional $V$, prove $\dim\bigcap U_i\ge\sum\dim U_i-(k-1)n$.
    \ii Prove or disprove: $\dim(U_1+U_2+U_3)=\sum\dim U_i-\sum_{i<j}\dim(U_i\cap U_j)+\dim(U_1\cap U_2\cap U_3)$.
    \ii Let $\dim V=n$ and $\dim U=m$.
    Prove $\mathcal E=\{T\in\LL(V):T(U)\subseteq U\}$ is a subspace of $\LL(V)$ and compute $\dim\mathcal E$.
  \end{enumerate}
\end{problem}

\begin{problem}[$\star$ Null spaces of compositions]
  Let $T\colon U\to V$ and $S\colon V\to W$ be linear with $U,V$ finite-dimensional.
  \begin{enumerate}[(a)]
    \ii Show $\nul(ST)=T^{-1}(\nul S)$, and deduce
    $\dim\nul(ST)=\dim\nul T+\dim(\nul S\cap\range T)$.
    \ii Prove $\dim\nul(ST)\le\dim\nul S+\dim\nul T$, with equality if and only if $\nul S\subseteq\range T$.
    \ii Prove $\rank(ST)\ge\rank S+\rank T-\dim V$.
    Determine exactly when $\rank(ST)=\rank T$.
    \ii Are there $S,T\in\LL(\RR^3)$ with $\rank S=\rank T=2$ and $ST=0$?
    Are there $S,T\in\LL(\FF^2)$ with $ST\ne0$ and $TS=0$?
  \end{enumerate}
\end{problem}

\begin{problem}[$\star$ Prescribing null spaces and ranges]
  \leavevmode
  \begin{enumerate}[(a)]
    \ii If $T\in\LL(V)$, $V$ finite-dimensional, and $T^2=0$, prove $\rank T\le\frac12\dim V$.
    \ii Prove some $T\in\LL(\FF^n)$ has $\nul T=\range T$ if and only if $n$ is even.
    Write one down for $n=6$.
    \ii Let $V,W$ be finite-dimensional and $U\subseteq V$.
    Prove some $T\in\LL(V,W)$ has $\nul T=U$ if and only if $\dim U\ge\dim V-\dim W$.
    \ii Prove no linear map $\FF^5\to\FF^2$ has null space $\{x: x_1=3x_2,\ x_3=x_4=x_5\}$.
    \ii If $\nul T\cap\range T=\{0\}$ for $T\in\LL(V)$, $V$ finite-dimensional, prove $V=\nul T\oplus\range T$.
  \end{enumerate}
\end{problem}

\begin{problem}[Matrices of polynomial maps]
  \leavevmode
  \begin{enumerate}[(a)]
    \ii Define $\Delta\colon\PP_3(\RR)\to\PP_2(\RR)$ by $(\Delta p)(x)=p(x+1)-p(x)$.
    Explain why $\deg\Delta p\le2$.
    Find $\MM(\Delta)$ in the monomial bases, and find $\nul\Delta$ and $\range\Delta$.
    \ii Find $p$ with $\Delta p=x^2$ and $p(0)=0$. Use it to compute $\sum_{k=0}^{n-1}k^2$.
    \ii Find $\dim\{p\in\PP_3(\RR):\Delta p\in\spn(1,x^2)\}$ using the preimage formula.
    \ii Let $T=I+D$ on $\PP_3(\RR)$, $Dp=p'$.
    Prove $T$ is an isomorphism without a matrix, show $T^{-1}=I-D+D^2-D^3$, and solve $p+p'=x^3$.
    \ii Is $I+D$ invertible on $\PP(\RR)$? Compare with $D$ on $\PP(\RR)$.
  \end{enumerate}
\end{problem}

\begin{problem}[$\star$ Interpolation with derivatives]
  \leavevmode
  \begin{enumerate}[(a)]
    \ii For $a\in\RR$ let $T_a\colon\PP_2(\RR)\to\RR^3$, $T_a(p)=(p(0),p(1),p'(a))$.
    Determine exactly which $a$ make $T_a$ an isomorphism.
    For the others, find $\nul T_a$ and $\range T_a$, and explain using the mean value theorem.
    \ii Prove, before computing, that exactly one $p\in\PP_3(\RR)$ has
    $p(0)=1$, $p'(0)=0$, $p(1)=2$, $p'(1)=-1$. Then find it.
    \ii Prove or disprove: for distinct $x_0,x_1$ and any $a_0,a_1,b_0,b_1$, exactly one $p\in\PP_3(\RR)$ has $p(x_i)=a_i$, $p'(x_i)=b_i$.
  \end{enumerate}
\end{problem}

\begin{problem}[$\star$ Invertibility and one-sided inverses]
  \leavevmode
  \begin{enumerate}[(a)]
    \ii Over $\FF_5$, for which $a$ is $A_a=\begin{psmallmatrix}a&2\\3&a\end{psmallmatrix}$ invertible?
    Find $A_2^{-1}$.
    \ii Over $\CC$, is $B=\begin{psmallmatrix}1+i&2\\1&1-i\end{psmallmatrix}$ invertible?
    If not, find a nonzero vector in $\nul B$ and describe $\range B$.
    \ii Count the invertible $2\times2$ matrices over $\Fp$.
    \ii Let $A=\begin{psmallmatrix}1&2&3\\0&1&1\end{psmallmatrix}\in\RR^{2,3}$.
    Find all $B$ with $AB=I_2$. Prove no $C$ has $CA=I_3$.
    \ii Prove $A\in\FF^{m,n}$ has a right inverse if and only if $\rank A=m$, and a left inverse if and only if $\rank A=n$.
  \end{enumerate}
\end{problem}

\begin{problem}[$\star$ Structure theorems for isomorphisms]
  All spaces are finite-dimensional.
  \begin{enumerate}[(a)]
    \ii Let $T_1,T_2\in\LL(V,W)$.
    Prove $\nul T_1=\nul T_2$ if and only if $T_1=ST_2$ for some invertible $S\in\LL(W)$.
    \ii Prove $\range T_1=\range T_2$ if and only if $T_1=T_2S$ for some invertible $S\in\LL(V)$.
    \ii If $ST=TS$ for every $S\in\LL(V)$, prove $T=\lambda I$.
    \ii Over $\RR$, suppose $T^2-3T+2I=0$.
    Prove $T$ is invertible and write $T^{-1}$ as a polynomial in $T$.
    If $V\ne\{0\}$, prove $T-I$ or $T-2I$ is not invertible.
    What goes wrong over $\FF_2$?
    \ii If $\dim V\ge2$, prove the non-invertible operators on $V$ do not form a subspace.
  \end{enumerate}
\end{problem}

\begin{problem}[$\star$ Products, affine sets, extensions]
  \leavevmode
  \begin{enumerate}[(a)]
    \ii Prove two translates $v+U$, $w+U$ of a subspace $U$ are equal or disjoint.
    \ii Prove $\{x:Tx=b\}$ is empty or equal to $x_0+\nul T$ for any solution $x_0$.
    Use this to list all solutions over $\FF_2$ of $x+y+z+w=1$, $x+z=0$.
    \ii Prove $V\times W=(V\times\{0\})\oplus(\{0\}\times W)$ and $V\cong V\times\{0\}$.
    Reprove $\dim(V\times W)=\dim V+\dim W$.
    \ii Is $\PP_2(\RR)\times\RR^2\cong\RR^{1,5}$? Is $\LL(\RR^2,\RR^3)\cong\RR^2\times\RR^4$?
    \ii Let $U\subseteq V$, $V$ finite-dimensional, and $S\colon U\to V$ injective.
    Prove $S$ extends to an invertible $T\in\LL(V)$.
    Show the hypothesis ``injective'' cannot be dropped.
  \end{enumerate}
\end{problem}

\begin{problem}[Quotient spaces]
  \leavevmode
  \begin{enumerate}[(a)]
    \ii Let $U=\spn((1,1,1))\subseteq\RR^3$.
    Decide whether $(1,0,0)+U=(0,1,-1)+U$ and whether $(2,1,1)+U=(1,0,0)+U$.
    Find a basis of $\RR^3/U$, and write $(3,5,7)+U$ in that basis.
    \ii Prove: if $u_1,\dots,u_m$ is a basis of $U$ and $u_1,\dots,u_m,v_1,\dots,v_k$ is a basis of $V$, then $v_1+U,\dots,v_k+U$ is a basis of $V/U$.
    Conclude $\dim V/U=\dim V-\dim U$ without using rank--nullity.
    \ii On $\RR^2/\spn(e_2)$, show that $f\big((x,y)+U\big)=x$ is a well-defined isomorphism onto $\RR$.
    Show that the same formula is not well defined on $\RR^2/\spn(e_1)$.
    \ii Let $T\in\LL(V,W)$ and $U\subseteq V$ a subspace.
    Prove $\tilde T(v+U)=Tv$ defines a map $V/U\to W$ if and only if $U\subseteq\nul T$.
    In that case, find $\nul\tilde T$.
    \ii Let $T\colon\PP_3(\RR)\to\RR^2$, $T(p)=(p(0),p(1))$.
    Find $\nul T$ and a basis of $\PP_3(\RR)/\nul T$, and describe the isomorphism $\PP_3(\RR)/\nul T\to\RR^2$ from the first isomorphism theorem.
  \end{enumerate}
\end{problem}

\section{True or false}
Give a one-line reason or counterexample. Spaces are finite-dimensional unless stated otherwise.
\begin{enumerate}
  \ii Over $\FF_3$, every nonempty subset of $\FF_3^n$ closed under addition is a subspace.
  \ii Over $\RR$, every nonempty subset of $\RR^n$ closed under addition is a subspace.
  \ii If $U\cup W$ is a subspace, then $U\subseteq W$ or $W\subseteq U$.
  \ii $(1,1,0),(0,1,1),(1,0,1)$ is a basis of $\FF_2^3$.
  \ii There is an injective linear map $\PP_4(\RR)\to\RR^4$.
  \ii If $ST$ is injective, then $S$ is injective.
  \ii If $T^2=0$, then $I-T$ is invertible.
  \ii On an infinite-dimensional space, every injective operator is surjective.
  \ii If $\dim V=5$ and $T\in\LL(V,\FF^3)$, then $\dim\nul T\ge2$.
  \ii For every subspace $U\subseteq V$ there is $T\in\LL(V)$ with $\nul T=\range T=U$.
  \ii If $v_1,\dots,v_n$ is a basis and $T$ is surjective, then $Tv_1,\dots,Tv_n$ is a basis of $W$.
  \ii $\LL(\FF^2,\FF^3)\cong\FF^6$.
  \ii The non-invertible $2\times2$ matrices over $\FF_2$ form a subspace.
  \ii If $S,T\in\LL(V)$ and $ST=I$, then $TS=I$.
  \ii For $A\in\FF^{2,3}$, some $B$ has $AB=I_2$ if and only if the columns of $A$ span $\FF^2$.
\end{enumerate}

\newpage
\section{Solutions}\label{sec:sol}

\subsection*{Techniques}
\textbf{T1.}
(a) $T=I+D+D^2$ on $\PP_n$. If $p\ne0$ has degree $d$, then $p'$ and $p''$ have smaller degree, so $Tp$ has degree $d$ and is nonzero. $T$ is injective, hence an isomorphism: exactly one $p$ for each $q$.
(b) If $\big((x^2+1)p\big)''=0$, then $(x^2+1)p$ has degree $\le1$; but for $p\ne0$ it has degree $\ge2$. So the map is injective on $\PP_n$, hence surjective. For $p\mapsto(x^2+1)p''$: constants are killed, so it is not injective, so by T1 it is not surjective either.
(c) $T(p)=0$ means $p$ has $n+1$ roots and degree $\le n$, so $p=0$. Injective, equal dimensions, isomorphism. Surjectivity says every list of values $(b_0,\dots,b_n)$ is hit.

\textbf{T2.}
(a) $I,T,T^2,\dots,T^{n^2}$ are $n^2+1$ vectors in $\LL(V)$, which has dimension $n^2$. So $\sum_{i=0}^{n^2}a_iT^i=0$ with not all $a_i=0$; $p(x)=\sum a_ix^i$ works.
(b) Take a basis $v_1,\dots,v_n$ of $V$. The $n$ vectors $Tv_1,\dots,Tv_n$ lie in $W$ with $n>\dim W$, so some nontrivial $\sum a_iTv_i=0$. Then $v=\sum a_iv_i\ne0$ by independence, and $Tv=0$.

\textbf{T3.}
(a) Row 2 with $S=T$: $\dim\nul(T^2)=\dim\nul T+\dim(\nul T\cap\range T)\le2\dim\nul T$. Row 3 with $S=T$: $\rank(T^2)=\rank T-\dim(\nul T\cap\range T)\ge\rank T-\dim\nul T=2\rank T-\dim V$.
(b) $T|_X$ has null space $X\cap\nul T=\{0\}$, so it is injective, so it is an isomorphism onto its range $T(X)$. Then $\dim X=\dim T(X)\le\dim\range T$.
(c) This is row 2 with $S=T$. $\nul T\subseteq\nul(T^2)$ always, so equality holds iff the dimensions agree iff $\dim(\nul T\cap\range T)=0$.
(d) $\dim T(U)=\dim U-\dim(U\cap\nul T)\ge\dim U-\dim\nul T$. Equality when $\nul T\subseteq U$, e.g.\ $U=V$; strict when $U\cap\nul T=\{0\}$ but $\nul T\ne\{0\}$, e.g.\ $T=N$ on $\RR^2$ and $U=\spn(e_2)$.

\textbf{T4.}
(a) If all four values vanish, $p$ has double roots at $0$ and $1$, so $x^2(x-1)^2$ divides $p$; degree $4>3$ forces $p=0$. Injective, so an isomorphism.
(b) Isomorphism iff $a\ne\tfrac12$ (Problem 8(a)). At $a=\tfrac12$ the data is $p(0),p(1),p'(\tfrac12)$: a derivative at $\tfrac12$ without $p(\tfrac12)$, so the orders at $\tfrac12$ do not start from $0$, and T4 does not apply. And indeed $p'(\tfrac12)=p(1)-p(0)$ for every quadratic.
(c) $L_i(a_k)=\delta_{ik}$. If $\sum c_iL_i=0$, evaluate at $a_k$: $c_k=0$. So $n+1$ independent vectors in $\PP_n$, a basis. The interpolating polynomial is $p=\sum b_iL_i$, since $p(a_k)=b_k$.

\textbf{T5.}
(a) $x^k\mapsto2^kx^k$: diagonal $2^k\ne0$. $x^k\mapsto(x+1)^k=x^k+\dots$: diagonal all $1$. Both invertible.
(b) $\Delta(x^k)=(x+1)^k-x^k$ has degree exactly $k-1$ with leading coefficient $k$, for $k\ge1$; $\Delta(1)=0$. So $\range\Delta\subseteq\PP_{n-1}$ (each $\Delta(x^k)$ has degree $\le n-1$) and it contains a polynomial of each degree $0,\dots,n-1$, hence (distinct degrees are independent) has dimension $n$, so $\range\Delta=\PP_{n-1}$. Then $\dim\nul\Delta=(n+1)-n=1$, and the constants are in it.
(c) $T(x^k)=kx^k-x^k=(k-1)x^k$. The matrix is diagonal with entries $-1,0,1,2,\dots,n-1$. $\nul T=\spn(x)$; $\range T=\spn(1,x^2,x^3,\dots,x^n)$.

\textbf{T6.}
(a) On $\PP_3$, $D^4=0$, so $I-D+D^2-D^3=(I+D)^{-1}$. The equation says $(I+D)^{-1}p=x^3$, so $p=(I+D)x^3=x^3+3x^2$. Check: $p-p'+p''-p'''=(x^3+3x^2)-(3x^2+6x)+(6x+6)-6=x^3$.
(b) $N(x^k)=(k-1)x^k$, so $I-N$ has diagonal $2,1,0$ in the monomial basis: not invertible; it kills $x^2$. T6 does not apply because $N$ is not nilpotent: $Nx^2=x^2$, so $N^kx^2=x^2$ for all $k$.

\textbf{T7.}
(a) See Problem 10(d).
(b) $T\cdot T^2=I$, so $T^{-1}=T^2$. $(T-I)(T^2+T+I)=T^3-I=0$ with $T-I\ne0$; by T8 (the factors commute), $T^2+T+I$ is not injective.
(c) $T=0$: $0+0=0$, not invertible. $T=I$: $I+I=2I=0$ in $\FF_2$, invertible.

\textbf{T8.}
(a) Yes. $T$: $e_3\mapsto e_1$, $e_4\mapsto e_2$, $e_1,e_2\mapsto0$; $S$: $e_1,e_2\mapsto0$, $e_3\mapsto e_3$, $e_4\mapsto e_4$. Then $\range T=\spn(e_1,e_2)=\nul S$, so $ST=0$. (In $\RR^3$ it is impossible since $2+2>3$.)
(b) $e_2\mapsto e_1$, $e_4\mapsto e_3$, $e_1,e_3\mapsto0$ has rank $2$ and $T^2=0$. No larger rank: $T\cdot T=0$ gives $2\rank T\le4$.
(c) $P(P-I)=P^2-P=0$ with $P-I\ne0$, so $P$ is not injective. $(I-P)P=0$ with $P\ne0$, so $I-P$ is not injective.

\textbf{T9.}
(a) Extend $v_1,\dots,v_k$ to a basis $v_1,\dots,v_n$ of $V$. Define $Tv_i=w_i$ for $i\le k$ and $Tv_i=0$ for $i>k$ (LADR 3.4). Without independence: $v_1=v_2=v$ with $w_1\ne w_2$ is impossible, since $Tv$ has one value.
(b) Let $u_1,\dots,u_m$ be a basis of $U$, extended by $w_1,\dots,w_k$ to a basis of $V$.
$\range T=U$: define $Tw_i=u_i$ for $i\le\min(k,m)$, the remaining $w_i\mapsto0$, and $u_i\mapsto u_i$; the range is $\spn(u_1,\dots,u_m)=U$.
$\nul T=U$: define $Tu_i=0$ and $Tw_j=w_j$; then $T(\sum a_iu_i+\sum b_jw_j)=\sum b_jw_j$, which is $0$ only when all $b_j=0$.
Both: rank--nullity forces $\dim U=\dim V-\dim U$. If $k=m$, define $Tu_i=0$ and $Tw_i=u_i$: the range is $U$ and the null space contains $U$ with dimension $\dim V-m=m$, so it equals $U$.
(c) For independent $v,w$: $T(v+w)=a_{v+w}(v+w)$ and also $=a_vv+a_ww$; comparing coefficients, $a_v=a_{v+w}=a_w$. For $w=cv$ with $c\ne0$: $a_ww=T(cv)=ca_vv=a_vw$, so $a_w=a_v$. All the scalars agree; call it $\lambda$.

\textbf{T10.}
(a) $J(A)=A^T$ satisfies $J^2=I$, so T10 gives $A=\tfrac12(A+A^T)+\tfrac12(A-A^T)$ with the two parts symmetric and skew, and the sum is direct. Dimensions: a symmetric matrix is determined by its entries on and above the diagonal, $\tfrac12n(n+1)$ of them; a skew matrix has zero diagonal and is determined by the entries above it, $\tfrac12n(n-1)$. Check: they add to $n^2$.
(b) $J(p)(x)=p(-x)$. Even part $\tfrac12(p(x)+p(-x))$, odd part $\tfrac12(p(x)-p(-x))$. Even polynomials in $\PP_n$: monomials $1,x^2,\dots$, count $\lfloor n/2\rfloor+1$; odd: $\lceil n/2\rceil$.
(c) $(I-P)^2=I-2P+P^2=I-P$. $v\in\nul(I-P)$ iff $Pv=v$ iff $v\in\range P$ (if $v=Pw$ then $Pv=P^2w=v$). Apply the same to $I-P$ for the other statement.
(d) $N(x,y)=(y,0)$: $\nul N=\range N=\spn(e_1)$. $N^2=0\ne N$, so $N$ is not idempotent.

\textbf{T11.}
(a) $U=\spn((1,1))$. $x$: value $1$ on $(1,1)$, not well defined. $x-y$: value $0$ on $(1,1)$, well defined. $x+y$: value $2$, not well defined. $x^2-y^2$: not linear, so test directly: $(1,0)+U=(2,1)+U$ but the formula gives $1$ and $3$; not well defined. (Note $x^2-y^2=(x-y)(x+y)$; the $x+y$ factor is what breaks it.)
(b) Well defined by T11 since $f(U)=0$. $\bar f(v+U)=0$ iff $f(v)=0$ iff $v\in\nul f$. So $\bar f$ is injective iff every $v\in\nul f$ has $v+U=0+U$, i.e.\ $\nul f\subseteq U$. With $U\subseteq\nul f$ given, this says $U=\nul f$.
(c) Each $v$ has a unique decomposition $u+w$, so $P(v)=u$ is a well-defined function. Linear: the decomposition of $\lambda v+v'$ is $(\lambda u+u')+(\lambda w+w')$. $P(u)=u$ for $u\in U$, so $P^2=P$. $\range P=U$ and $\nul P=\{v:u=0\}=W$.

\textbf{T12.}
(a) $(8-1)(8-2)(8-4)=168$.
(b) Ordered independent pairs in $\FF_2^3$: $7\cdot6=42$. Each $2$-dimensional subspace has $3$ nonzero vectors and $3\cdot2=6$ ordered bases. So $42/6=7$ subspaces.
(c) Rank $2$, $3$ unknowns: $3^{1}=3$ solutions. $y=x-2=x+1$, $z=1-x-y=-2x=x$: $(0,1,0),(1,2,1),(2,0,2)$.
(d) The evaluation map $\PP_4(\FF_2)\to\{\text{functions }\FF_2\to\FF_2\}$ has target of dimension $2$ and is surjective (constants and $x$), so its null space has dimension $5-2=3$: $2^3=8$ polynomials, e.g.\ the span of $x^2+x$, $x^3+x$, $x^4+x$.

\textbf{T13.}
(a) (i) False: $N$. (ii) False: shifts. (iii) True (T1). (iv) False: three lines.
(b) $e_3\mapsto e_2\mapsto e_1\mapsto0$, $e_4\mapsto0$. $T^2e_3=e_1\ne0$, $T^3=0$. $\rank T=2$, $\rank T^2=1$.

\subsection*{Hammond's practice midterm (full-credit answers)}
Answers to the practice midterm distributed in class (4 questions, 40 points); the questions are not reproduced here.
\medskip\noindent\textbf{Q1(a) (6).}
$\dim U=2$ since $u_1,u_2$ are independent (neither is a multiple of the other). Likewise $\dim W=2$.
Suppose $v\in U\cap W$: $v=a u_1+b u_2=c w_1+d w_2$, so $(a,a+b,b,0)=(d,0,c,c-d)$.
Fourth coordinates: $c=d$. Second: $a=-b$. First: $a=d$. Third: $b=c$. Then $b=c=d=a=-b$, so $b=0$ and everything is $0$. Hence $U\cap W=\{0\}$, so $U+W$ is direct and $\dim(U\oplus W)=4=\dim\RR^4$, giving $U\oplus W=\RR^4$.

\noindent\textit{What costs points:} forgetting to argue $\dim U=\dim W=2$; showing $U\cap W=\{0\}$ without the dimension count that gives $U+W=\RR^4$; or trying instead to show $u_1,u_2,w_1,w_2$ is a basis of $\RR^4$ and asserting independence without proof. The count replaces that argument entirely.

\medskip\noindent\textbf{Q1(b) (2).}
Solve $(2,1,3,0)=a u_1+b u_2+c w_1+d w_2=(a+d,\,a+b,\,b+c,\,c-d)$: $c=d$, $a+d=2$, $a+b=1$, $b+c=3$. From the last two, $c-a=2$, and with $a+c=2$ we get $c=2,a=0$, then $b=1,d=2$.
So $u=u_2=(0,1,1,0)$ and $w=2w_1+2w_2=(2,0,2,0)$.

\medskip\noindent\textbf{Q2(a) (4).}
Suppose $\sum a_iTv_i=0$. By linearity $T(\sum a_iv_i)=0$, so $\sum a_iv_i\in\nul T=\{0\}$ since $T$ is injective. Independence of $v_1,\dots,v_n$ gives all $a_i=0$.

\medskip\noindent\textbf{Q2(b) (6).}
By (a), $Tv_1,\dots,Tv_n$ is independent in $W$; extend it to a basis $Tv_1,\dots,Tv_n,w_1,\dots,w_k$ of $W$ (possible since $W$ is finite-dimensional).
A linear map is determined by its values on a basis, so define $S\colon W\to V$ by $S(Tv_i)=v_i$ and $S(w_j)=0$.
Then $STv_i=v_i$ for each $i$, so $ST$ and $I_V$ agree on the basis $v_1,\dots,v_n$, hence $ST=I_V$.

\noindent\textit{What costs points:} defining $S$ only on $\range T$ without saying what it does on the rest of $W$, or saying ``let $S=T^{-1}$'' when $T$ is not surjective.

\medskip\noindent\textbf{Q3(a) (6).}
Apply $T$ to each vector of $B$ and read off coefficients in $1,x,x^2$:
$T(1,0,0)=1+x^2$, $T(1,1,0)=2+x+x^2$, $T(1,1,1)=2+2x+2x^2$. So
\[ \MM(T)=\begin{pmatrix}1&2&2\\0&1&2\\1&1&2\end{pmatrix}. \]

\medskip\noindent\textbf{Q3(b) (2).}
$\dim\RR^3=\dim\PP_2=3$, so it is enough to show $T$ is injective. If $T(a,b,c)=0$, then $a+b=b+c=a+c=0$; adding gives $a+b+c=0$, so $a=b=c=0$. By LADR 3.65, $T$ is an isomorphism.

\medskip\noindent\textbf{Q4(a) (4).}
If $v\in T^{-1}(U)$, then $Tv\in U$ by definition, and $Tv\in\range T$, so $S(v)\in U\cap\range T$.
$S$ is linear because $T$ is: $S(\lambda v+w)=T(\lambda v+w)=\lambda Tv+Tw$.
Surjective: if $y\in U\cap\range T$, write $y=Tv$; then $Tv\in U$, so $v\in T^{-1}(U)$ and $S(v)=y$.

\medskip\noindent\textbf{Q4(b) (2).}
$v\in\nul S$ means $v\in T^{-1}(U)$ and $Tv=0$. Since $0\in U$, every $v$ with $Tv=0$ is in $T^{-1}(U)$. So $\nul S=\nul T$.

\medskip\noindent\textbf{Q4(c) (4).}
The fundamental theorem applied to $S$ on the finite-dimensional space $T^{-1}(U)$:
$\dim T^{-1}(U)=\dim\nul S+\dim\range S=\dim\nul T+\dim(U\cap\range T)$, using (a) for the range and (b) for the null space.

\medskip\noindent\textbf{Q4(d) (4).}
$\nul D=\spn(1)$, dimension $1$. $\range D=\PP_3(\RR)$ (every polynomial of degree $\le3$ is a derivative), so $U\cap\range D=U$, dimension $2$. By (c), $\dim D^{-1}(U)=1+2=3$.
$D^{-1}(U)$ is the set of $p$ with $p'\in\spn(1,x^2)$; for $p=a+bx+cx^2+dx^3+ex^4$, $p'=b+2cx+3dx^2+4ex^3$, so $c=e=0$. Basis: $1,x,x^3$.

\subsection*{Practice Midterm A}
\textbf{1a.} $U$ is the null space of the linear map $p\mapsto p(0)$, which is surjective ($1\mapsto1$), so $\dim U=3-1=2$ (basis $x,x^2$). $\dim W=1$.
If $c(1+x)\in U$, then its value at $0$ is $c=0$, so $U\cap W=\{0\}$.
Then $\dim(U\oplus W)=2+1=3$, so $U\oplus W=\PP_2(\RR)$.

\textbf{1b.} $w=2(1+x)$, $u=x+x^2$.

\textbf{2a.} If $\sum a_jv_j=0$, apply $T$: $\sum a_jw_j=0$, so all $a_j=0$ since the $w_j$ form a basis.

\textbf{2b.} Define $S$ on the basis by $Sw_j=v_j$. Then $TSw_j=Tv_j=w_j$, so $TS$ and $I_W$ agree on a basis, hence $TS=I_W$.

\textbf{2c.} $T(a,b)=a$. Both $S_1(t)=(t,0)$ and $S_2(t)=(t,t)$ satisfy $TS=I$.

\textbf{3a.} To write $(\alpha,\beta,\gamma)$ in the basis $C$: $a(1,0,0)+b(1,1,0)+c(1,1,1)=(a+b+c,\,b+c,\,c)$, so $c=\gamma$, $b=\beta-\gamma$, $a=\alpha-\beta$.
$T(1)=(1,1,1)\mapsto(0,0,1)$; $T(x)=(0,1,-1)\mapsto(-1,2,-1)$; $T(x^2)=(0,1,1)\mapsto(-1,0,1)$.
\[ \MM(T)=\begin{pmatrix}0&-1&-1\\0&2&0\\1&-1&1\end{pmatrix}. \]

\textbf{3b.} Dimensions are equal ($3$). If $T(p)=0$, then $p$ has roots $0,1,-1$ but degree $\le2$, so $p=0$. Injective, hence an isomorphism (LADR 3.65).

\textbf{3c.} $p(0)=0$ gives $p=bx+cx^2$; $p(1)=b+c=1$ and $p(-1)=-b+c=0$ give $b=c=\tfrac12$. So $p=\tfrac12(x+x^2)$.

\textbf{4a.} If $u\in\nul(ST)$, then $S(Tu)=0$, so $Tu\in\nul S$; and $Tu\in\range T$. $R$ is linear because $T$ is.
Surjective: if $y\in\nul S\cap\range T$, write $y=Tu$; then $STu=Sy=0$, so $u\in\nul(ST)$ and $R(u)=y$.

\textbf{4b.} $\nul R=\{u\in\nul(ST):Tu=0\}$. Every $u$ with $Tu=0$ has $STu=0$, so this is all of $\nul T$.

\textbf{4c.} Rank--nullity for $R$: $\dim\nul(ST)=\dim\nul R+\dim\range R=\dim\nul T+\dim(\nul S\cap\range T)$.

\textbf{4d.} $\nul T=\spn(1)$, dimension $1$. $\range T=\PP_2(\RR)$.
$\nul S=\{p:p(0)=p(1)=0\}$; its intersection with $\PP_2(\RR)$ is $\spn\big(x(x-1)\big)$, dimension $1$ (a quadratic with roots $0,1$ is a multiple of $x(x-1)$).
So $\dim\nul(ST)=1+1=2$.
Directly: $STp=(p'(0),p'(1))$. For $p=a+bx+cx^2+dx^3$, $p'=b+2cx+3dx^2$, so $b=0$ and $2c+3d=0$. Basis: $1,\ 2x^3-3x^2$.

\subsection*{Practice Midterm B}
\textbf{1a.} Every $A$ equals $\tfrac12(A+A^T)+\tfrac12(A-A^T)$, the first symmetric and the second skew, so $V=U+W$.
If $A\in U\cap W$, then $A=A^T=-A$, so $2A=0$ and $A=0$. So the sum is direct.
$\dim U=3$ (basis $E_{11},E_{22},E_{12}+E_{21}$), $\dim W=1$ (basis $E_{12}-E_{21}$).

\textbf{1b.} $u=\begin{psmallmatrix}1&5/2\\5/2&4\end{psmallmatrix}$, $w=\begin{psmallmatrix}0&-1/2\\1/2&0\end{psmallmatrix}$.

\textbf{1c.} $\nul T=U$. $(A-A^T)^T=-(A-A^T)$, so $\range T\subseteq W$; and for $B\in W$, $T(\tfrac12B)=\tfrac12(B-B^T)=B$. So $\range T=W$. (Check: $3+1=4$.)

\textbf{2a.} $\Gamma\big(\lambda(u,w)+(u',w')\big)=\lambda u+u'+\lambda w+w'=\lambda\Gamma(u,w)+\Gamma(u',w')$. The outputs are exactly the vectors $u+w$, which is $U+W$.

\textbf{2b.} $\Gamma(u,w)=0$ means $w=-u$. Then $w\in W$ and $w=-u\in U$, so $u\in U\cap W$ and $(u,w)=(u,-u)$. Conversely every $(u,-u)$ with $u\in U\cap W$ is in $U\times W$ and maps to $0$.
The map $u\mapsto(u,-u)$ is linear, injective (the first coordinate recovers $u$), and surjective onto $\nul\Gamma$ by the description just proved.

\textbf{2c.} Rank--nullity for $\Gamma$: $\dim U+\dim W=\dim(U\times W)=\dim\nul\Gamma+\dim\range\Gamma=\dim(U\cap W)+\dim(U+W)$.

\textbf{2d.} If $U\cap W=\{0\}$, then $\nul\Gamma=\{0\}$, so $\Gamma$ is injective. It is surjective onto $U+W=U\oplus W$. So $\Gamma$ is an isomorphism $U\times W\to U\oplus W$.

\textbf{3a.} $T(1)=1$, $T(x)=1+x$, $T(x^2)=1+2x+x^2$:
$\MM(T)=\begin{psmallmatrix}1&1&1\\0&1&2\\0&0&1\end{psmallmatrix}$.

\textbf{3b.} $T(1)=1$; $T(x)=x+1=1\cdot1+1\cdot x$; $T\big(x(x-1)\big)=(x+1)x=x^2+x=x(x-1)+2x$.
$\MM(T)=\begin{psmallmatrix}1&1&0\\0&1&2\\0&0&1\end{psmallmatrix}$.

\textbf{3c.} If $p(x+1)=0$ for all $x$, then $p=0$, so $T$ is injective; equal dimensions give an isomorphism. $(T^{-1}p)(x)=p(x-1)$: both compositions give $p(x)$ back.

\textbf{4a.} If $w=Tv$, then $Tw=T^2v=0$.

\textbf{4b.} By (a) and rank--nullity, $\dim V=\dim\nul T+\dim\range T\ge2\dim\range T$.

\textbf{4c.} $(I-T)(I+T)=I-T^2=I$ and $(I+T)(I-T)=I$. So $(I-T)^{-1}=I+T$.

\textbf{4d.} $T^2p=p''''=0$ for $\deg p\le3$. For $p=a+bx+cx^2+dx^3$, $p''=2c+6dx$, so $\range T=\PP_1(\RR)$ and $\nul T=\{c=d=0\}=\PP_1(\RR)$. Here range $=$ null and $2\le\tfrac12\cdot4$.
$p=(I+T)x^3=x^3+6x$. Check: $p''=6x$ and $p-p''=x^3$.

\subsection*{Practice Midterm C}
\textbf{1a.} $U$ is the null space of the linear map $E\colon\PP_3(\RR)\to\RR$, $E(p)=p(2)$, so it is a subspace. $E$ is surjective ($E(1)=1$), so $\dim U=4-1=3$.

\textbf{1b.} $x-2,\ (x-2)^2,\ (x-2)^3$. Each vanishes at $2$. They are independent because their degrees $1,2,3$ are different (compare leading terms in a relation). Three independent vectors in a $3$-dimensional space form a basis.

\textbf{1c.} $W=\spn(1)$. $1\notin U$, so $U\cap W=\{0\}$, and $3+1=4$.

\textbf{2a.} $\range(ST)=\{S(Tu)\}=\{Sy:y\in\range T\}=S(\range T)$, and every $Sy$ is in $\range S$.

\textbf{2b.} $\range(ST)\subseteq\range S$ gives the first bound. Let $S'$ be $S$ restricted to $\range T$; its range is $S(\range T)=\range(ST)$, so by rank--nullity $\dim\range(ST)\le\dim\range T$.

\textbf{2c.} With $S'$ as in (b): $\nul S'=\nul S\cap\range T\subseteq\nul S$, so $\dim\nul S'\le\dim\nul S=\dim V-\dim\range S$.
Rank--nullity for $S'$: $\dim\range(ST)=\dim\range T-\dim\nul S'\ge\dim\range T-\dim V+\dim\range S$.

\textbf{3a.} $T(1)=(1,1,0)$, $T(x)=(0,1,1)$, $T(x^2)=(0,1,2)$:
$\MM(T)=\begin{psmallmatrix}1&0&0\\1&1&1\\0&1&2\end{psmallmatrix}$.

\textbf{3b.} Equal dimensions. If $T(p)=0$ with $p=a+bx+cx^2$: $a=0$, $b+c=0$, $b+2c=0$, so $c=0$ then $b=0$. Injective, hence an isomorphism.

\textbf{3c.} $p=1+bx+cx^2$ with $1+b+c=0$ and $b+2c=2$. Subtracting, $c=3$, $b=-4$: $p=1-4x+3x^2$.

\textbf{4a.} $v+\nul T=v'+\nul T$ means $v-v'\in\nul T$, so $Tv-Tv'=T(v-v')=0$.

\textbf{4b.} Linear: $\tilde T\big((v+N)+(w+N)\big)=\tilde T(v+w+N)=T(v+w)=Tv+Tw$, and similarly for scalars.
Injective: if $\tilde T(v+N)=0$, then $Tv=0$, so $v\in N$ and $v+N=0+N$.
Surjective: every element of $\range T$ is some $Tv=\tilde T(v+N)$.

\textbf{4c.} Isomorphic finite-dimensional spaces have the same dimension.

\textbf{4d.} $T(a+bx+cx^2)=b+c$, so $\nul T=\{b=-c\}$ with basis $1,\ x-x^2$.
$\dim\PP_2/\nul T=3-2=1$; since $T(x)=1\ne0$, $x\notin\nul T$, so $x+\nul T$ is a basis.
$x-x^2\in\nul T$, so $x+\nul T=x^2+\nul T$. Yes.

\subsection*{Problem 1}
\begin{enumerate}[(a)]
  \ii ($\Rightarrow$) If $x\in f^{-1}(f(U))$ then $f(x)=f(u)$ for some $u\in U$, so $x=u\in U$.
  ($\Leftarrow$) If $f(x_1)=f(x_2)$, take $U=\{x_1\}$: then $x_2\in f^{-1}(f(U))=U$, so $x_2=x_1$.
  \ii ($\Rightarrow$) For $y\in V$ pick $x$ with $f(x)=y$; then $x\in f^{-1}(V)$.
  ($\Leftarrow$) Take $V=Y$: $Y=f(f^{-1}(Y))=f(X)$.
  \ii The inclusion is immediate.
  If $f$ is injective and $y=f(a)=f(b)$ with $a\in A$, $b\in B$, then $a=b\in A\cap B$.
  If $f(x_1)=f(x_2)$ with $x_1\ne x_2$, then $A=\{x_1\}$, $B=\{x_2\}$ give $f(A\cap B)=\varnothing\ne f(A)\cap f(B)$.
  \ii $v\in T^{-1}(T(U))\iff Tv=Tu$ for some $u\in U\iff v-u\in\nul T$ for some $u\in U\iff v\in U+\nul T$.
  And $w\in T(T^{-1}(X))\iff w=Tv$ with $Tv\in X\iff w\in\range T\cap X$.
  \ii $T|_U$ has null space $U\cap\nul T$ and range $T(U)$; apply rank--nullity.
  By (d), $T^{-1}(T(U))=U\iff\nul T\subseteq U$.
  This holds for $U=\{0\}$ if and only if $\nul T=\{0\}$.
\end{enumerate}

\subsection*{Problem 2}
\begin{enumerate}[(a)]
  \ii $2\cdot3=0$ with $2,3\ne0$, impossible in a field.
  \ii $ax=ay\implies a(x-y)=0\implies x=y$.
  An injective self-map of a finite set is bijective, and it fixes $0$, so it permutes the nonzero $x_1,\dots,x_{q-1}$.
  Then $a^{q-1}\prod x_i=\prod(ax_i)=\prod x_i$, and $\prod x_i\ne0$ cancels: $a^{q-1}=1$.
  So $a^q=a$ for all $a$ (including $0$).
  \ii The indicator functions $\delta_c$ form a basis, so $\dim\SF=p$; $\dim\PP_n=n+1$.
  \begin{itemize}
    \ii Injective if and only if $n\le p-1$: a nonzero polynomial of degree $\le n<p$ cannot vanish at all $p$ points.
    If $n\ge p$, then $x^p-x\in\nul E$ by (b).
    \ii Surjective if and only if $n\ge p-1$: at $n=p-1$, $E$ is injective between $p$-dimensional spaces, so surjective; larger $n$ only enlarges the range.
    For $n<p-1$ the domain is too small.
    \ii Isomorphism if and only if $n=p-1$.
  \end{itemize}
  \ii $(x+1)^4=(x^2+1)^2=x^4+1$, which equals $x+1$ as a function since $c^4=c$ on $\FF_2$.
  Ex.~34 needs $n+1$ distinct points in $\FF$; with $|\FF|<n+1$ they do not exist, and distinct polynomials (e.g.\ $x^2$ and $x$ over $\FF_2$) define the same function.
\end{enumerate}

\subsection*{Problem 3}
\begin{enumerate}[(a)]
  \ii False: the vectors sum to $0$.
  \ii If $\sum c_jw_j=0$, the coefficient of $v_j$ is $c_j+c_{j-1}$ (indices mod $n$), so $c_j=(-1)^{j-1}c_1$, and wrapping around gives $c_1=(-1)^nc_1$.
  \begin{itemize}
    \ii $n$ even: $c=(1,-1,\dots,-1)$ works, so the list is dependent over every field.
    \ii $n$ odd: $2c_1=0$. Independent if and only if $2\ne0$ in $\FF$; dependent over $\FF_2$.
  \end{itemize}
  When independent, $w_1-w_2+w_3-\dots+w_n=2v_1$, so $v_1=\tfrac12(w_1-w_2+\dots+w_n)$.
  \ii Take $\sum a_i(v_i+w)=0$, not all $a_i=0$.
  Then $\sum a_iv_i=-(\sum a_i)w$.
  If $s=\sum a_i=0$, independence forces all $a_i=0$, contradiction.
  So $w=-s^{-1}\sum a_iv_i$.
  \ii True: $\spn(v_1,\dots,v_n)$ has dimension $n$, and a spanning list of length $n$ in it is a basis.
\end{enumerate}

\subsection*{Problem 4}
\begin{enumerate}[(a)]
  \ii $\dim(U\cap W)\ge7+7-10=4$, then $\dim(U\cap W\cap X)\ge4+7-10=1$.
  Equality: in $\FF^{10}$, $U=\{x_1=x_2=x_3=0\}$, $W=\{x_4=x_5=x_6=0\}$, $X=\{x_7=x_8=x_9=0\}$ meet in $\spn(e_{10})$.
  \ii Induct with $\dim(I\cap U_k)\ge\dim I+\dim U_k-n$ where $I=U_1\cap\dots\cap U_{k-1}$.
  \ii False: lines spanned by $e_1,e_2,e_1+e_2$ in $\RR^2$ give $2$ on the left and $3$ on the right.
  \ii Closure: $(\lambda S+T)u=\lambda Su+Tu\in U$.
  Extend a basis $u_1,\dots,u_m$ of $U$ to a basis of $V$.
  $T(U)\subseteq U$ if and only if rows $m+1,\dots,n$ of columns $1,\dots,m$ of $\MM(T)$ vanish.
  Since $\MM$ is an isomorphism, $\dim\mathcal E=n^2-m(n-m)$.
\end{enumerate}

\subsection*{Problem 5}
\begin{enumerate}[(a)]
  \ii $u\in\nul(ST)\iff Tu\in\nul S$. Apply the preimage formula (HW3 P1) with $U=\nul S$.
  \ii $\nul S\cap\range T\subseteq\nul S$ gives the inequality.
  Equality if and only if the two nested subspaces are equal, if and only if $\nul S\subseteq\range T$.
  \ii $\rank(ST)=\dim U-\dim\nul(ST)\ge\dim U-\dim\nul T-\dim\nul S=\rank T+\rank S-\dim V$.
  Exactly: $\rank(ST)=\rank T-\dim(\nul S\cap\range T)$, so $\rank(ST)=\rank T$ if and only if $\nul S\cap\range T=\{0\}$.
  \ii No: $\rank(ST)\ge2+2-3=1$.
  Yes: $S=\begin{psmallmatrix}1&0\\0&0\end{psmallmatrix}$, $T=\begin{psmallmatrix}0&1\\0&0\end{psmallmatrix}$ give $ST=T\ne0$ and $TS=0$.
\end{enumerate}

\subsection*{Problem 6}
\begin{enumerate}[(a)]
  \ii $Tw=T^2v=0$ for $w=Tv$, so $\range T\subseteq\nul T$, and $\dim V\ge2\rank T$.
  \ii ($\Rightarrow$) $n=2\rank T$.
  ($\Leftarrow$) For $n=2k$ set $Te_j=0$ and $Te_{k+j}=e_j$ ($j\le k$).
  Then $\range T=\spn(e_1,\dots,e_k)\subseteq\nul T$ and both have dimension $k$.
  For $n=6$: $\MM(T)=\begin{psmallmatrix}0&I_3\\0&0\end{psmallmatrix}$.
  \ii ($\Rightarrow$) $\dim V-\dim U=\rank T\le\dim W$.
  ($\Leftarrow$) Extend a basis $u_1,\dots,u_m$ of $U$ by $v_1,\dots,v_k$, where $k\le\dim W$.
  Set $Tu_i=0$ and $Tv_i=w_i$ for independent $w_1,\dots,w_k$.
  If $T(\sum a_iu_i+\sum b_iv_i)=\sum b_iw_i=0$, all $b_i=0$, so $\nul T=U$.
  \ii The subspace $\{(3s,s,t,t,t)\}$ has dimension $2<5-2$.
  \ii The sum is direct and has dimension $\dim\nul T+\dim\range T=\dim V$.
\end{enumerate}

\subsection*{Problem 7}
\begin{enumerate}[(a)]
  \ii The $x^k$ terms of $(x+1)^k-x^k$ cancel.
  $\Delta1=0$, $\Delta x=1$, $\Delta x^2=2x+1$, $\Delta x^3=3x^2+3x+1$, so
  \[ \MM(\Delta)=\begin{pmatrix}0&1&1&1\\0&0&2&3\\0&0&0&3\end{pmatrix}. \]
  Rank $3$, so $\range\Delta=\PP_2(\RR)$ and $\nul\Delta=\spn(1)$.
  \ii Solving gives $c_3=\tfrac13$, $c_2=-\tfrac12$, $c_1=\tfrac16$:
  $p=\tfrac13x^3-\tfrac12x^2+\tfrac16x$.
  Telescoping, $\sum_{k=0}^{n-1}k^2=p(n)-p(0)=\tfrac{n(n-1)(2n-1)}{6}$.
  \ii $1+\dim(\spn(1,x^2)\cap\PP_2)=1+2=3$.
  \ii If $p\ne0$ has degree $d$, then $p+p'$ has degree $d$, so $T$ is injective, hence an isomorphism.
  $D^4=0$ on $\PP_3$, so $(I+D)(I-D+D^2-D^3)=I-D^4=I$.
  Then $p=(I-D+D^2-D^3)x^3=x^3-3x^2+6x-6$.
  \ii Yes: $Sp=\sum_{k\ge0}(-D)^kp$ is a finite sum for each $p$, and $(I+D)S=S(I+D)=I$.
  $D$ itself is surjective but not injective on $\PP(\RR)$.
\end{enumerate}

\subsection*{Problem 8}
\begin{enumerate}[(a)]
  \ii For $p=c_0+c_1x+c_2x^2$, $T_ap=0$ means $c_0=0$, $c_1+c_2=0$, $c_1+2ac_2=0$, so $(2a-1)c_2=0$.
  For $a\ne\tfrac12$, $T_a$ is injective, hence an isomorphism.
  For $a=\tfrac12$: $\nul T_a=\spn(x-x^2)$, and since $p(1)-p(0)=c_1+c_2=p'(\tfrac12)$,
  $\range T_a=\{(y_0,y_1,z):z=y_1-y_0\}$.
  For a quadratic, the secant slope on $[0,1]$ equals the derivative at the midpoint exactly, so $p'(\tfrac12)$ cannot be prescribed.
  \ii $Sp=(p(0),p'(0),p(1),p'(1))$ on $\PP_3\to\RR^4$.
  If $Sp=0$ then $p=c_2x^2+c_3x^3$ with $c_2+c_3=0=2c_2+3c_3$, so $p=0$.
  Hence $S$ is an isomorphism.
  Solving: $p=1+4x^2-3x^3$.
  \ii True: a $p\in\PP_3$ with double roots at $x_0,x_1$ is divisible by the degree-$4$ polynomial $(x-x_0)^2(x-x_1)^2$, so $p=0$. So $p\mapsto(p(x_0),p'(x_0),p(x_1),p'(x_1))$ is injective $\PP_3\to\RR^4$, equal dimensions, hence an isomorphism: exactly one $p$.
\end{enumerate}

\subsection*{Problem 9}
\begin{enumerate}[(a)]
  \ii $\det A_a=a^2-6=(a-1)(a+1)$ in $\FF_5$, so invertible if and only if $a\in\{0,2,3\}$.
  For $a=2$, $\det=3$ and $3^{-1}=2$:
  $A_2^{-1}=2\begin{psmallmatrix}2&-2\\-3&2\end{psmallmatrix}=\begin{psmallmatrix}4&1\\4&4\end{psmallmatrix}$.
  \ii $\det B=(1+i)(1-i)-2=0$.
  $(2,-1-i)^T\in\nul B$.
  Column 2 is $(1-i)$ times column 1, so $\range B=\spn\big((1+i,1)^T\big)$.
  \ii $(p^2-1)(p^2-p)$: any nonzero first column, then a second column that is not a multiple of it.
  \ii $\nul A=\spn\big((-1,-1,1)^T\big)$; particular solutions $(1,0,0)^T$ for $e_1$ and $(-2,1,0)^T$ for $e_2$. So
  \[ B=\begin{pmatrix}1-s&-2-t\\-s&1-t\\s&t\end{pmatrix},\quad s,t\in\RR. \]
  No left inverse: $\rank(CA)\le\rank A\le2<3$.
  \ii $AB=I_m\implies m\le\rank A\le m$.
  Conversely, if $A$ is surjective, choose $b_j$ with $Ab_j=e_j$.
  $CA=I_n\implies n\le\rank A\le n$.
  Conversely, if $A$ is injective, extend $Ae_1,\dots,Ae_n$ to a basis of $\FF^m$ and let $C$ send $Ae_j\mapsto e_j$ and the rest to $0$.
\end{enumerate}

\subsection*{Problem 10}
\begin{enumerate}[(a)]
  \ii ($\Leftarrow$) $S$ is injective, so $ST_2v=0\iff T_2v=0$.
  ($\Rightarrow$) Take a basis $u_1,\dots,u_m$ of the common null space, extended by $v_1,\dots,v_k$.
  Both $T_1v_1,\dots,T_1v_k$ and $T_2v_1,\dots,T_2v_k$ are independent.
  Extend them to bases $T_2v_i,x_j$ and $T_1v_i,y_j$ of $W$ (same length).
  Define $S(T_2v_i)=T_1v_i$ and $Sx_j=y_j$: basis to basis, so $S$ is invertible, and $ST_2=T_1$ on $u_i$ and $v_i$.
  \ii ($\Leftarrow$) $\range(T_2S)=T_2(V)$ since $S$ is surjective.
  ($\Rightarrow$) Equal ranks give equal nullities $m$.
  Take an adapted basis $u_i,v_i$ for $T_1$, choose $v_i'$ with $T_2v_i'=T_1v_i$, and a basis $u_i'$ of $\nul T_2$.
  Then $u_i',v_i'$ is independent (apply $T_2$ to a relation) and has length $\dim V$.
  Set $Su_i=u_i'$ and $Sv_i=v_i'$.
  \ii If $v,Tv$ were independent, extend to a basis and let $Sv=v$, $S(Tv)=0$: then $STv=0\ne Tv=TSv$.
  So $Tv=a_vv$ for every $v$.
  For independent $v,w$: $a_{v+w}(v+w)=a_vv+a_ww$ gives $a_v=a_w$.
  For $w=cv$: $a_ww=T(cv)=a_vw$.
  So $T=\lambda I$.
  \ii $T(3I-T)=2I$, so $T^{-1}=\tfrac12(3I-T)$.
  $(T-I)(T-2I)=0$ is not invertible when $V\ne\{0\}$, so the factors cannot both be.
  Over $\FF_2$ the equation is $T^2-T=0$; dividing by $2$ is impossible and $T=0$ satisfies it.
  \ii With a basis $v_1,\dots,v_n$, let $P_1v_1=v_1$, $P_1v_j=0$ for $j\ge2$, and $P_2=I-P_1$.
  Neither is injective, but $P_1+P_2=I$.
\end{enumerate}

\subsection*{Problem 11}
\begin{enumerate}[(a)]
  \ii If $v+u_1=w+u_2$, then $v-w\in U$, so $v+U=w+U$ (LADR 3.101).
  \ii $Tx=b\iff T(x-x_0)=0$.
  Over $\FF_2$: $z=x$, then $y+w=1$. With $x_0=(0,0,0,1)$ and $\nul=\spn((1,0,1,0),(0,1,0,1))$, the solutions are
  $(0,0,0,1)$, $(1,0,1,1)$, $(0,1,0,0)$, $(1,1,1,0)$.
  \ii $(v,w)=(v,0)+(0,w)$ and the intersection is $\{0\}$.
  $v\mapsto(v,0)$ is a linear bijection onto $V\times\{0\}$.
  Dimensions of a direct sum of two subspaces add.
  \ii Both yes: dimensions $5=5$ and $6=6$.
  \ii This is HW4 P7; the proof is in the review (Isomorphisms). Counterexample without injectivity: $S=0$ on $U=V=\FF$; any extension is $0$, not invertible.
\end{enumerate}

\subsection*{Problem 12}
\begin{enumerate}[(a)]
  \ii $(1,0,0)-(0,1,-1)=(1,-1,1)\notin U$, so the translates are different (and disjoint).
  $(2,1,1)-(1,0,0)=(1,1,1)\in U$, so they are equal.
  $(1,1,1),e_1,e_2$ is a basis of $\RR^3$, so $e_1+U,e_2+U$ is a basis of $\RR^3/U$ by (b).
  $(3,5,7)=7(1,1,1)-4e_1-2e_2$, so $(3,5,7)+U=-4(e_1+U)-2(e_2+U)$.
  \ii Spanning: write $v=\sum a_iu_i+\sum b_jv_j$; then $v+U=\sum b_j(v_j+U)$ because $\sum a_iu_i\in U$.
  Independence: if $\sum b_j(v_j+U)=0+U$, then $\sum b_jv_j\in U$, so $\sum b_jv_j=\sum a_iu_i$ for some $a_i$.
  Independence of the full basis forces all $b_j=0$.
  So $\dim V/U=k=\dim V-\dim U$.
  \ii If $(x,y)+U=(x',y')+U$, then $(x-x',y-y')\in\spn(e_2)$, so $x=x'$: well defined.
  It is linear, surjective ($x\mapsto(x,0)+U$), and injective ($x=0$ means $(0,y)\in U$).
  On $\RR^2/\spn(e_1)$: $(0,0)+U=(1,0)+U$ but the formula gives $0$ and $1$.
  \ii ($\Leftarrow$) If $v-v'\in U\subseteq\nul T$, then $Tv=Tv'$.
  Linearity follows from linearity of $T$.
  ($\Rightarrow$) If $u\in U$ has $Tu\ne0$, then $0+U=u+U$ would have to map to both $0$ and $Tu$.
  $\nul\tilde T=\{v+U:v\in\nul T\}$, which is $(\nul T)/U$.
  In particular $\tilde T$ is injective if and only if $U=\nul T$.
  \ii $T$ is surjective ($p=a+(b-a)x$ hits $(a,b)$), so $\dim\nul T=4-2=2$ and $\nul T=\spn\big(x(x-1),\,x^2(x-1)\big)$.
  $1,x,x(x-1),x^2(x-1)$ is a basis of $\PP_3(\RR)$, so $1+\nul T,\ x+\nul T$ is a basis of the quotient.
  The isomorphism is $p+\nul T\mapsto(p(0),p(1))$, sending $1+\nul T\mapsto(1,1)$ and $x+\nul T\mapsto(0,1)$.
  Two polynomials are equal in the quotient if and only if they agree at $0$ and $1$.
\end{enumerate}

\subsection*{True or false}
\begin{enumerate}
  \ii True: $2v=v+v$ and $0=3v$ lie in the set, and those are all the scalar multiples.
  \ii False: $\ZZ^n$.
  \ii True: if $u\in U\setminus W$, $w\in W\setminus U$, then $u+w$ lies in neither.
  \ii False: the sum is $0$. (True over $\RR$.)
  \ii False: $5>4$.
  \ii False: $T(t)=(t,0)$, $S(x,y)=x$ give $ST=I$. ($T$ must be injective.)
  \ii True: $(I-T)^{-1}=I+T$.
  \ii False: the right shift.
  \ii True: $5-\rank T\ge5-3$.
  \ii False: requires $\dim U=\frac12\dim V$.
  \ii False: spans but need not be independent; $(x,y)\mapsto x$.
  \ii True: $\dim=6$.
  \ii False: $E_{11}+E_{22}=I$.
  \ii True in finite dimensions; false for the shifts on $\FF^\infty$.
  \ii True: a right inverse exists if and only if $\rank A=2$.
\end{enumerate}

\end{document}
